% Luc Destombe --- licence CC BY-NC-SA 4.0
% https://creativecommons.org/licenses/by-nc-sa/4.0/deed.fr
% Fichier autonome : compiler avec pdflatex, deux fois (signets, nombre de pages).
\def\lyceeversion{1.0}
\documentclass[11pt,a4paper]{article}
\makeatletter
\newif\ifeleve
\elevefalse

\newif\iflycee@palatino
\lycee@palatinotrue

\RequirePackage[utf8]{inputenc}
\RequirePackage[T1]{fontenc}
\RequirePackage[french]{babel}
\RequirePackage{etoolbox}

\newcommand{\ShorthandsOff}{\@ifundefined{shorthandoff}{}{\shorthandoff{;:!?}}}
\newcommand{\ShorthandsOn}{\@ifundefined{shorthandon}{}{\shorthandon{;:!?}}}
\ShorthandsOff
\AtBeginDocument{\ShorthandsOn}
\AtBeginEnvironment{tikzpicture}{\ShorthandsOff}

\RequirePackage{geometry}
\RequirePackage{fancyhdr}
\RequirePackage{lastpage}
\RequirePackage{multicol}
\RequirePackage{array}
\RequirePackage{enumitem}
\RequirePackage{xcolor}
\RequirePackage{needspace}

\RequirePackage{amsmath,amssymb}
\RequirePackage{mathpazo}
\DeclareMathSizes{10.95}{10.95}{8}{5.5}
\begingroup\catcode`\_=\active \catcode`\^=\active
\gdef_#1{\sb{\scriptscriptstyle #1}}
\gdef^#1{\sp{\scriptscriptstyle\lycee@moinsepais #1}}
\endgroup
\newcommand\lycee@moins{\mathbin{%
  \setbox\z@\hbox{$\m@th\scriptscriptstyle\lycee@moinsorig$}%
  \dimen@=.55\wd\z@ \advance\dimen@-.5pt
  \hbox to.75\wd\z@{\hss$\m@th\scriptscriptstyle\vcenter{\hbox{%
    \tikz\draw[line width=.5pt,line cap=round](0,0)--(\the\dimen@,0);}}$\hss}}}
\begingroup\lccode`\~=`\- \lowercase{\endgroup\let~}\lycee@moins
\newcommand\lycee@moinsepais{\mathcode`\-="8000 }
\AtBeginDocument{\mathchardef\lycee@moinsorig=\mathcode`\-\relax}
\AtBeginDocument{\catcode`\_=12 \mathcode`\_="8000
  \catcode`\^=12 \mathcode`\^="8000 }
\newcommand\lycee@exposants{%
  \fontdimen13\textfont2=.48\fontdimen6\textfont2
  \fontdimen14\textfont2=.48\fontdimen6\textfont2
  \fontdimen15\textfont2=.48\fontdimen6\textfont2
  \fontdimen13\scriptfont2=.48\fontdimen6\scriptfont2
  \fontdimen14\scriptfont2=.48\fontdimen6\scriptfont2
  \fontdimen15\scriptfont2=.48\fontdimen6\scriptfont2}
\every@math@size\expandafter{\the\every@math@size\lycee@exposants}
\newlength{\retraitcalculs}
\setlength{\retraitcalculs}{2em}

\RequirePackage{tikz}
\usetikzlibrary{arrows.meta,calc,babel,shapes.misc}
\RequirePackage[most]{tcolorbox}
\tcbuselibrary{skins,breakable}

\definecolor{coulDef}{HTML}{1F4E79}
\definecolor{coulProp}{HTML}{7030A0}
\definecolor{coulRem}{HTML}{C55A11}
\definecolor{coulEx}{HTML}{2E7D32}
\definecolor{coulExo}{HTML}{B03A2E}
\definecolor{coulPart}{HTML}{1F4E79}
\definecolor{coulMeth}{HTML}{1B6E4F}
\definecolor{coulCourbe}{HTML}{0B5ED7}
\definecolor{coulCourbeB}{HTML}{C00000}
\definecolor{coulCourbeC}{HTML}{2E7D32}
\definecolor{coulGrille}{HTML}{8C8C8C}
\definecolor{coulRep}{HTML}{1B6E4F}
\definecolor{coulBar}{HTML}{2E7D32}

\newcommand{\R}{\mathbb{R}}
\newcommand{\Rs}{\mathbb{R}^{*}}
\renewcommand{\le}{\leqslant}
\renewcommand{\ge}{\geqslant}
\renewcommand{\approx}{\simeq}
\newcommand{\vect}[1]{\overrightarrow{#1}}
\newcommand{\norme}[1]{\left\lVert #1 \right\rVert}
\newcommand{\intervalle}[2]{\left[#1\,;#2\right]}
\RequirePackage{cancel}
\renewcommand{\CancelColor}{\color{coulExo}}
\newcommand{\fois}[1]{\times{\color{coulExo}#1}}
\newcommand{\barre}[1]{\cancel{#1}}

\tikzset{grille/.style={coulGrille,line width=0.4pt}}
\newcommand{\quadrillage}[4]{\draw[grille,step=1] (#1,#2) grid (#3,#4);}
\tikzset{
  croix/.style={cross out,draw,line width=0.6pt,inner sep=0pt,minimum size=4pt},
  croixdroite/.style={croix,rotate=45,line width=0.8pt},
}
\newcommand{\pt}[3]{\path (#1) node[croix]{} node[#2,font=\small,inner sep=2.5pt]{$#3$};}

\newcommand{\lycee@repere}[4]{%
  \draw[grille,step=1] (#1,#2) grid (#3,#4);
  \draw[-{Stealth[length=2mm]},thick] (#1,0)--({#3+0.4},0) node[below left=-1pt]{$x$};
  \draw[-{Stealth[length=2mm]},thick] (0,#2)--(0,{#4+0.4}) node[below left=-1pt]{$y$};
  \node[below left,font=\scriptsize,inner sep=1.5pt] at (0,0) {$0$};}
\newcommand{\lycee@gradx}[1]{\foreach \k/\l in {#1}{%
  \draw (\k,0.07)--(\k,-0.07) node[below,font=\scriptsize,inner sep=1.5pt]{$\l$};}}
\newcommand{\lycee@grady}[1]{\foreach \k/\l in {#1}{%
  \draw (0.07,\k)--(-0.07,\k) node[left,font=\scriptsize,inner sep=1.5pt]{$\l$};}}
\newcommand{\lycee@intff}[2]{\left[#1\,;#2\right]}
\newcommand{\lycee@intfo}[2]{\left[#1\,;#2\right[}
\newcommand{\lycee@intof}[2]{\left]#1\,;#2\right]}
\newcommand{\lycee@intoo}[2]{\left]#1\,;#2\right[}
\newcommand{\lycee@ens}[1]{\left\{#1\right\}}
\AtBeginDocument{%
  \providecommand{\repere}{\lycee@repere}%
  \providecommand{\gradx}{\lycee@gradx}%
  \providecommand{\grady}{\lycee@grady}%
  \providecommand{\Cf}{\mathcal{C}\sb{\scriptscriptstyle f}}%
  \providecommand{\Cg}{\mathcal{C}\sb{\scriptscriptstyle g}}%
  \providecommand{\intff}{\lycee@intff}%
  \providecommand{\intfo}{\lycee@intfo}%
  \providecommand{\intof}{\lycee@intof}%
  \providecommand{\intoo}{\lycee@intoo}%
  \providecommand{\ens}{\lycee@ens}}

\newlist{questions}{enumerate}{3}
\setlist[questions,1]{label=\arabic*\textdegree),leftmargin=*,itemsep=2pt,topsep=3pt}
\setlist[questions,2]{label=\alph*),leftmargin=*,itemsep=1pt,topsep=2pt}
\setlist[questions,3]{label=\roman*),leftmargin=*,itemsep=1pt,topsep=1pt}
\setlist[itemize]{label=\textbullet,leftmargin=*,itemsep=2pt,topsep=3pt}

\newcommand{\besoinplace}[1]{\par\penalty\@M\begingroup
  \dimen@=#1\relax\dimen@ii=\pagegoal\advance\dimen@ii-\pagetotal
  \ifdim\dimen@>\dimen@ii\ifdim\dimen@ii>\z@\vfil\fi\break\fi\endgroup}

\newcommand{\lycee@pied}{\thepage/\pageref{LastPage}}
\newcommand{\entete}[2]{%
  \pagestyle{fancy}\fancyhf{}%
  \fancyhead[L]{\footnotesize\color{coulPart}\textbf{#1}}%
  \fancyhead[R]{\footnotesize\color{coulPart}\textbf{#2}}%
  \fancyfoot[C]{\footnotesize\lycee@pied}%
  \renewcommand{\headrulewidth}{0.4pt}%
  \renewcommand{\headrule}{\hbox to\headwidth{%
    \color{coulPart!40}\leaders\hrule height \headrulewidth\hfill}}}

\RequirePackage{ccicons}
\newcommand{\lycee@licence}{{\scriptsize\color{gray}%
  \copyright\ Luc Destombe --- \ccbyncsa\ CC BY-NC-SA 4.0}}
\newif\iflycee@licence \lycee@licencetrue
\newcommand{\sanslicence}{\lycee@licencefalse}
\AtBeginDocument{\iflycee@licence\AddToHookNext{shipout/foreground}{%
  \put(0,0){\rlap{\kern\dimexpr1in+\hoffset+\oddsidemargin+\textwidth\relax
    \llap{\raisebox{-\dimexpr1in+\voffset+\topmargin+\headheight+\headsep
      +\textheight+\footskip\relax}{\lycee@licence}}}}}\fi}

\newcommand{\lycee@site}{faire-des-maths.fr}
\newcommand{\lycee@version}{{\scriptsize\color{gray}\lycee@site
  \ifdefined\lyceeversion\ --- Version~\lyceeversion\fi}}
\AtBeginDocument{\iflycee@licence\AddToHookNext{shipout/foreground}{%
  \put(0,0){\rlap{\kern\dimexpr1in+\hoffset+\oddsidemargin\relax
    \raisebox{-\dimexpr1in+\voffset+\topmargin+\headheight+\headsep
      +\textheight+\footskip\relax}{\lycee@version}}}}\fi}

\newcommand{\formulesserrees}{%
  \setlength{\abovedisplayskip}{3pt}\setlength{\belowdisplayskip}{3pt}%
  \setlength{\abovedisplayshortskip}{2pt}\setlength{\belowdisplayshortskip}{3pt}}

\setlength{\parindent}{0pt}

\RequirePackage[implicit=false,hidelinks,pdfencoding=auto,psdextra,bookmarksopen,
  bookmarksopenlevel=1,pdfcreator={},pdfproducer={}]{hyperref}
\RequirePackage{bookmark}
\hypersetup{pdfauthor={Luc Destombe},pdfkeywords={CC BY-NC-SA 4.0}}
\newcounter{lycee@signet}
\newcommand{\signet}[3][]{\stepcounter{lycee@signet}%
  \edef\lycee@dest{\if\relax\detokenize{#1}\relax
    signet.\arabic{lycee@signet}\else#1\fi}%
  \hypertarget{\lycee@dest}{}\bookmark[level=#2,dest=\lycee@dest]{#3}}
\pdfstringdefDisableCommands{%
  \def\R{R}\def\vect#1{#1}\def\Cf{Cf}\def\Cg{Cg}\def\fois#1{×#1}%
  \def\barre#1{#1}\def\intervalle#1#2{[#1;#2]}\def\intff#1#2{[#1;#2]}%
  \def\textsuperscript#1{#1}\def\dfrac#1#2{#1/#2}\def\frac#1#2{#1/#2}%
  \def\vec#1{#1⃗}}%
\geometry{top=2cm,bottom=1cm,left=1cm,right=1cm,headheight=14pt,footskip=0.6cm}

\RequirePackage{pgfplots}
\pgfplotsset{compat=1.18}
\RequirePackage{tkz-tab}

\newcounter{partiecours}
\renewcommand{\thepartiecours}{\Roman{partiecours}}
\newcommand{\partiecours}[1]{%
  \refstepcounter{partiecours}%
  \Needspace*{8\baselineskip}%
  \bigskip
  \noindent\begin{tcolorbox}[enhanced,colback=coulPart,colframe=coulPart,
      boxrule=0pt,arc=2pt,left=8pt,right=8pt,top=4pt,bottom=4pt,
      before skip=6pt,after skip=10pt]
    \color{white}\bfseries\large \thepartiecours{} --- #1%
    \signet[partie-\arabic{partiecours}]{1}{\thepartiecours{} --- #1}
  \end{tcolorbox}%
  \addcontentsline{toc}{section}{\thepartiecours{} --- #1}%
}

\newtcolorbox{boitecours}[3][]{%
  enhanced,breakable,
  colback=white,colframe=#2,coltitle=white,
  fonttitle=\bfseries,
  attach boxed title to top left={xshift=8pt,yshift=-\tcboxedtitleheight/2},
  boxed title style={colback=#2,arc=2pt,boxrule=0pt},
  boxrule=0.9pt,arc=2pt,
  left=8pt,right=8pt,top=8pt,bottom=6pt,
  title={#3},#1}

\newcounter{methode}
\newenvironment{definitionbox}[1][Définition]%
  {\begin{boitecours}[phantom={\signet{2}{#1}}]{coulDef}{#1}}{\end{boitecours}}
\newenvironment{proprietebox}[1][Propriété]%
  {\begin{boitecours}[phantom={\signet{2}{#1}}]{coulProp}{#1}}{\end{boitecours}}
\newenvironment{methodebox}[1][Méthode]{\stepcounter{methode}%
  \begin{boitecours}[phantom={\signet[methode-\arabic{methode}]{2}{#1}}]{coulMeth}{#1}}%
  {\end{boitecours}}

\newcommand{\etape}[1]{\par\smallskip\textbf{\color{coulMeth}#1}\par\nobreak\smallskip}
\newcommand{\enonce}[1]{\emph{#1}\par\smallskip}
\newcommand{\reponse}[1]{\par\smallskip
  {\footnotesize\itshape\color{black!55}Réponses : #1}}
\newcommand{\demo}[1]{\textbf{\color{coulMeth}Démonstration.} #1}
\newcommand{\afaire}[1]{%
  \par\smallskip
  \noindent{\small\color{coulExo}$\blacktriangleright$\ \textbf{À faire :}
    fiche d'exercices du chapitre, #1.}\par\smallskip}

\newenvironment{calculs}{\setlength{\jot}{5pt}\@fleqntrue\@mathmargin\retraitcalculs
  \setlength{\abovedisplayskip}{4pt}\setlength{\belowdisplayskip}{4pt}%
  \setlength{\abovedisplayshortskip}{4pt}\setlength{\belowdisplayshortskip}{4pt}%
  \csname alignat*\endcsname{2}}%
  {\csname endalignat*\endcsname}
\newcommand{\rem}[1]{\text{\small\itshape\color{black!60}#1}}
\newcommand{\explique}[2]{\par\smallskip\noindent
  \begin{minipage}[t]{0.57\linewidth}\raggedright #1\end{minipage}\hfill
  \begin{minipage}[t]{0.40\linewidth}\raggedright\leavevmode
    {\small\itshape\color{black!60}#2}\end{minipage}\par}
\newcommand{\cas}[1]{\par\smallskip\noindent\textbf{\color{coulExo}#1}\vadjust{\nobreak}}
\newcommand{\calc}[1]{\par\smallskip\textbf{\color{coulExo}#1}\par\nobreak}

\newtcolorbox{remarquebox}[1][Remarque]{%
  enhanced,breakable,colback=white,colframe=coulRem,
  boxrule=0pt,leftrule=2.5pt,arc=0pt,outer arc=0pt,
  left=8pt,right=6pt,top=5pt,bottom=5pt,
  before upper={\textbf{\color{coulRem}#1}\par\smallskip}}

\newtcolorbox{exemplebox}[1][Exemple]{%
  enhanced,breakable,colback=white,colframe=coulEx,
  boxrule=0pt,leftrule=2.5pt,arc=0pt,outer arc=0pt,
  left=8pt,right=6pt,top=5pt,bottom=5pt,
  before upper={\textbf{\color{coulEx}#1}\par\smallskip}}

\pgfplotsset{
  repere/.style={
    axis lines=middle,
    axis line style={-{Stealth[length=2.4mm]},thick},
    every axis x label/.style={at={(ticklabel* cs:1.0)},anchor=west},
    every axis y label/.style={at={(ticklabel* cs:1.0)},anchor=south},
    label style={font=\footnotesize},
    tick label style={font=\scriptsize},
    grid=both,
    major grid style={grille},
    minor tick num=0,
    tick align=inside,
    clip mode=individual,
  },
}
\tikzset{
  courbe/.style={coulCourbe,line width=1.1pt,smooth},
  courbeB/.style={coulCourbeB,line width=1.1pt,smooth},
  courbeC/.style={coulCourbeC,line width=1.1pt,smooth},
}

\newlist{etapes}{enumerate}{1}
\setlist[etapes]{label=\textbf{\arabic*.},leftmargin=*,itemsep=1pt,topsep=2pt}

\newlist{expressions}{enumerate}{1}
\setlist[expressions]{label={\Alph*\ $=$},leftmargin=*,itemsep=3pt,topsep=3pt}

\newcounter{exo}
\renewcommand{\theexo}{\ifnum\value{exo}<10 0\fi\arabic{exo}}
\newtcolorbox{exercice}[1][]{%
  enhanced,breakable,
  colback=white,colframe=coulExo,
  boxrule=0.9pt,arc=2pt,
  left=8pt,right=8pt,top=8pt,bottom=6pt,
  coltitle=white,fonttitle=\bfseries,
  attach boxed title to top left={xshift=8pt,yshift=-\tcboxedtitleheight/2},
  boxed title style={colback=coulExo,arc=2pt,boxrule=0pt},
  phantom={\signet{2}{Exercice \arabic{exo}}},
  title={Exercice \theexo\ifx\relax#1\relax\else{} --- #1\fi}}
\newenvironment{exo}[1][\relax]{\refstepcounter{exo}\begin{exercice}[#1]}{\end{exercice}}

  \newenvironment{correction}{\par}{\par}
  \newcommand{\autableau}{}
  \newcommand{\etapeexemples}{\etape{Exemples corrigés}}
  \newcommand{\etapetableau}[1]{\etape{#1}}
  \newenvironment{exempletableau}[1][Exemple]%
    {\begin{exemplebox}[#1]}{\end{exemplebox}}

\RequirePackage{comment}
\excludecomment{correctionavj}

\newcommand{\coursEntete}{}
\newcommand{\coursNiveau}{}
\newcommand{\coursTitre}{}
\newcommand{\coursSurTitre}{}
\newcommand{\coursSousTitre}{}

\newcommand{\enteteCours}[1]{\renewcommand{\coursEntete}{#1}}
\newcommand{\chapitrecours}[2]{\enteteCours{Chapitre #1 --- #2}}
\newcommand{\niveaucours}[1]{\renewcommand{\coursNiveau}{#1}}
\newcommand{\titrecours}[3][]{%
  \renewcommand{\coursTitre}{#2}%
  \renewcommand{\coursSurTitre}{#1}%
  \renewcommand{\coursSousTitre}{#3}}

\pagestyle{fancy}
\fancyhf{}
\fancyhead[L]{\footnotesize\color{coulPart}\textbf{\coursEntete}}
\fancyhead[R]{\footnotesize\color{coulPart}\textbf{\coursNiveau}}
\fancyfoot[C]{\footnotesize\thepage}
\renewcommand{\headrulewidth}{0.4pt}
\renewcommand{\headrule}{\hbox to\headwidth{\color{coulPart!40}\leaders\hrule height \headrulewidth\hfill}}

\setlength{\parskip}{2pt}

\AtBeginDocument{%
  \begin{center}
    \begin{tcolorbox}[enhanced,width=0.72\linewidth,colback=white,colframe=coulPart,
        boxrule=1.6pt,arc=3pt,halign=center,top=10pt,bottom=10pt]
      \ifx\coursSurTitre\empty
        {\Huge\bfseries\color{coulPart} \coursTitre}\\[4pt]
      \else
        {\Huge\bfseries\color{coulPart} \coursTitre}\\[3pt]
        {\Large\bfseries\color{coulPart} \coursSurTitre}\\[5pt]
      \fi
      {\normalsize \coursSousTitre}%
      
    \end{tcolorbox}
  \end{center}%
}

\usetikzlibrary{decorations.text}
\definecolor{lycee@metal}{HTML}{1F4E79}
\definecolor{lycee@manche}{HTML}{8A5A2B}
\definecolor{lycee@etincelle}{HTML}{F28C28}
\newcommand{\lycee@pluriel}[2]{#1 #2\ifnum#1>1 s\fi}
\newcommand{\lycee@sceau}{%
  \edef\lycee@bas{\lycee@pluriel{\lyceerelectures}{relecture} et
    \lycee@pluriel{\lyceerodages}{rodage}}%
  \begin{tikzpicture}[line join=round]
    \fill[coulPart!10,draw=coulPart,line width=.8pt] (0:2.2)
      \foreach \i in {1,...,71}{ -- (\i*5:{isodd(\i)?2.1:2.2}) } -- cycle;
    \draw[coulPart,line width=1.2pt] (0,0) circle (1.93);
    \draw[coulPart,line width=.4pt] (0,0) circle (1.85);
    \draw[coulPart,line width=.4pt] (0,0) circle (1.2);
    \path[decorate,decoration={text along path,text align=center,
        text={|\fontsize{9.5}{11}\selectfont\bfseries\color{coulPart}|TEST{\'E} AVEC MES {\'E}L{\`E}VES}}]
      (180:1.45) arc (180:0:1.45);
    \path[decorate,decoration={text along path,text align=center,
        text/.expanded={|\noexpand\fontsize{8.5}{10}\noexpand\selectfont\noexpand\color{coulPart}|\lycee@bas}}]
      (195:1.68) arc (195:345:1.68);
    \begin{scope}[x=.25mm,y=-.25mm]
      \begin{scope}[shift={(-32,-40)}]%
        \fill[lycee@metal] (23,31) rectangle (41,49);
        \draw[lycee@etincelle,line width=.5pt,line cap=round]
          (20,30)--(15,27) (19,34)--(14,35) (44,30)--(49,27) (45,34)--(50,35);
        \fill[lycee@manche,rounded corners=.3mm] (38,24.6) rectangle (60,27.6);
        \fill[lycee@metal,rounded corners=.4mm] (25,20) rectangle (39,31);
        \fill[lycee@etincelle,rounded corners=.2mm] (25,20) rectangle (28,31);
      \end{scope}
    \end{scope}
    \ifdefined\lyceesceauplusieurs\def\lycee@annee{Ann{\'e}es~:}\else
      \def\lycee@annee{Ann{\'e}e}\fi
    \path[decorate,decoration={text along path,text align=center,
        text/.expanded={|\noexpand\fontsize{8}{10}\noexpand\selectfont\noexpand\color{coulPart}|%
          \unexpanded\expandafter{\lycee@annee} \unexpanded\expandafter{\lyceesceauannees}}}]
      (95:1.08) arc (95:445:1.08);
  \end{tikzpicture}}
\AtEndDocument{\ifdefined\lyceesceauannees
  \par\besoinplace{5cm}\vspace{1.5em}%
  \begin{center}\lycee@sceau\end{center}\fi}
\makeatother

\chapitrecours{7}{Trigonométrie}
\niveaucours{1\textsuperscript{re} STI2D}
\titrecours{TRIGONOMÉTRIE}{Cours --- Première STI2D}

\begin{document}

\ShorthandsOff
\tikzset{
  lecture/.style={dashed,black!60,line width=0.6pt},
  sens/.style={-{Stealth[length=1.8mm]},line width=0.9pt,coulCourbeB},
  sensB/.style={-{Stealth[length=1.8mm]},line width=0.9pt,coulCourbeC},
}
\newcommand{\cerclenu}{%
  \draw[-{Stealth[length=2mm]},thick] (-1.3,0)--(1.3,0) node[below,font=\scriptsize,inner sep=2pt]{$x$};
  \draw[-{Stealth[length=2mm]},thick] (0,-1.3)--(0,1.3) node[right,font=\scriptsize,inner sep=2pt]{$y$};
  \draw[line width=0.9pt] (0,0) circle (1);
  \node[below left,font=\scriptsize,inner sep=1.5pt] at (0,0) {$O$};}
\newcommand{\cerclebase}{\cerclenu
  \node[below right,font=\scriptsize,inner sep=1.5pt] at (1,0) {$1$};
  \node[below left,font=\scriptsize,inner sep=1.5pt] at (-1,0) {$-1$};
  \node[above right,font=\scriptsize,inner sep=1.5pt] at (0,1) {$1$};
  \node[below right,font=\scriptsize,inner sep=1.5pt] at (0,-1) {$-1$};}
\newcommand{\rayons}[1]{\foreach \a in {#1}{%
  \draw[grille] (0,0)--({cos(\a)},{sin(\a)});}}
\newcommand{\ptc}[2]{%
  \path[coulCourbeB] ({cos(#1)},{sin(#1)}) node[croix]{};
  \node[font=\scriptsize,inner sep=1pt,align=center] at ({1.3*cos(#1)},{1.3*sin(#1)}) {#2};}
\newcommand{\verticale}[1]{\draw[lecture] ({cos(#1)},-1.15)--({cos(#1)},1.15);}
\newcommand{\horizontale}[1]{\draw[lecture] (-1.15,{sin(#1)})--(1.15,{sin(#1)});}
\newcommand{\lien}[2]{\draw[lecture] ({cos(#1)},{sin(#1)})--({cos(#2)},{sin(#2)});}
\newcommand{\lierx}[1]{\draw[lecture] ({cos(#1)},{sin(#1)})--({cos(#1)},0);}
\newcommand{\liery}[1]{\draw[lecture] ({cos(#1)},{sin(#1)})--(0,{sin(#1)});}
\newcommand{\valx}[2]{\pgfmathtruncatemacro{\q}{cos(#1)>0?1:0}%
  \draw[coulCourbeB,line width=1pt] ({cos(#1)},0.05)--({cos(#1)},-0.05);
  \ifnum\q=1 \node[below left,coulCourbeB,font=\scriptsize,inner sep=1pt,outer sep=1pt,fill=white] at ({cos(#1)},0) {#2};
  \else \node[below right,coulCourbeB,font=\scriptsize,inner sep=1pt,outer sep=1pt,fill=white] at ({cos(#1)},0) {#2};\fi}
\newcommand{\valy}[2]{%
  \draw[coulCourbeB,line width=1pt] (0.05,{sin(#1)})--(-0.05,{sin(#1)});
  \node[left,coulCourbeB,font=\scriptsize,inner sep=1pt,outer sep=1pt,fill=white] at (0,{sin(#1)}) {#2};}
\newcommand{\equacos}[4]{\begin{tikzpicture}[x=1.5cm,y=1.5cm]
  \cerclebase\verticale{#1}\valx{#1}{#4}\ptc{#1}{#2}\ptc{-#1}{#3}\end{tikzpicture}}
\newcommand{\equasin}[5]{\begin{tikzpicture}[x=1.5cm,y=1.5cm]
  \cerclebase\horizontale{#1}\valy{#1}{#5}\ptc{#1}{#3}\ptc{#2}{#4}\end{tikzpicture}}
\newcommand{\equapt}[4]{\begin{tikzpicture}[x=1.5cm,y=1.5cm]
  \cerclebase\verticale{#1}\horizontale{#1}\valx{#1}{#3}\valy{#1}{#4}\ptc{#1}{#2}\end{tikzpicture}}
\newcommand{\cerclept}[3]{\begin{tikzpicture}[x=1.5cm,y=1.5cm]
  \cerclebase#3\ptc{#1}{#2}\end{tikzpicture}}
\newcommand{\ptstable}[1]{\foreach \a in {-150,-135,-120,-90,-60,-45,-30,0,30,45,60,90,120,135,150,180}{%
  \path[coulCourbeB] ({\a*pi/180},{#1(\a)}) node[croix]{};}}
\newcommand{\arcsens}[4][sens]{%
  \draw[#1] ({#4*cos(#2)},{#4*sin(#2)}) arc (#2:#3:#4);}
\newcommand{\cerclemesure}[4]{%
  \begin{tikzpicture}[x=1.25cm,y=1.25cm]
    \rayons{#1}\cerclebase
    \arcsens{#3}{#2}{0.55}
    \ptc{#2}{#4}
  \end{tikzpicture}}
\newcommand{\figassociee}[3]{%
  \begin{tikzpicture}[x=1.25cm,y=1.25cm]
    #3\cerclebase
    \ptc{35}{$x$}
    \ptc{#1}{#2}
  \end{tikzpicture}}
\newcommand{\reperetrigo}[2]{%
  \foreach \k in {#1,...,#2}{\draw[grille] ({\k*pi/6},-1.5)--({\k*pi/6},1.5);}
  \foreach \j in {-3,...,3}{\draw[grille] ({#1*pi/6},{\j/2})--({#2*pi/6},{\j/2});}
  \draw[-{Stealth[length=2mm]},thick] ({#1*pi/6},0)--({#2*pi/6+0.4},0)
    node[below left=-1pt,font=\footnotesize]{$x$};
  \draw[-{Stealth[length=2mm]},thick] (0,-1.5)--(0,1.8)
    node[below left=-1pt,font=\footnotesize]{$y$};
  \node[above left,font=\scriptsize,inner sep=1pt,outer sep=1pt] at (0,1) {$1$};
  \foreach \j/\l in {0.5/\frac{1}{2},-0.5/-\frac{1}{2},-1/-1}{%
    \node[left,font=\scriptsize,inner sep=1pt,outer sep=1pt,fill=white] at (0,\j) {$\l$};}}
\newcommand{\gradpi}[2][below]{\foreach \k/\l in {#2}{%
  \draw ({\k*pi/6},0.05)--({\k*pi/6},-0.05);
  \node[#1,font=\scriptsize,inner sep=1pt,outer sep=2.5pt,fill=white] at ({\k*pi/6},0) {$\l$};}}
\newcommand{\gradcos}{%
  \gradpi{-12/-2\pi,-6/-\pi,6/\pi,12/2\pi}
  \gradpi[below left]{-9/-\frac{3\pi}{2},3/\frac{\pi}{2}}
  \gradpi[below right]{-3/-\frac{\pi}{2},9/\frac{3\pi}{2}}}
\newcommand{\gradsin}{%
  \gradpi{-9/-\frac{3\pi}{2},-3/-\frac{\pi}{2},3/\frac{\pi}{2},9/\frac{3\pi}{2}}
  \gradpi[below left]{-12/-2\pi,-6/-\pi,6/\pi}
  \gradpi[below right]{12/2\pi}}
\newcommand{\courbecos}[2]{\draw[coulCourbe,line width=1.1pt]
  plot[domain=#1*pi/6:#2*pi/6,samples=121] (\x,{cos(\x r)});}
\newcommand{\courbesin}[2]{\draw[coulCourbe,line width=1.1pt]
  plot[domain=#1*pi/6:#2*pi/6,samples=121] (\x,{sin(\x r)});}
\newcommand{\droiteh}[4]{\draw[coulCourbeB,line width=0.8pt]
  ({#1*pi/6},#3)--({#2*pi/6},#3) node[right,font=\footnotesize]{#4};
  \ifdim#3pt=0.5pt\node[left,font=\scriptsize,inner sep=1pt,outer sep=1pt,fill=white] at (0,0.5) {$\frac{1}{2}$};\fi
  \ifdim#3pt=-0.5pt\node[left,font=\scriptsize,inner sep=1pt,outer sep=1pt,fill=white] at (0,-0.5) {$-\frac{1}{2}$};\fi}
\newcommand{\lirex}[2]{\draw[lecture] ({#1*pi/6},#2)--({#1*pi/6},0);
  \path[coulCourbeB] ({#1*pi/6},#2) node[croix]{};}
\newcommand{\interaxe}[3][red]{\draw[#1,line width=1.8pt] ({#2*pi/6},0)--({#3*pi/6},0);}
\newcommand{\motif}[2]{\begin{scope}
  \clip (-5,-3.4) rectangle (5,3.4);
  \foreach \d in {#1}{\draw[coulCourbe,line width=1.1pt,shift={(\d,0)}] #2;}
  \end{scope}}
\newcommand{\reperet}[4]{%
  \draw[grille,step=1] (#1,#2) grid (#3,#4);
  \draw[-{Stealth[length=2mm]},thick] (#1,0)--({#3+0.4},0) node[below left=-1pt,font=\footnotesize]{$t$};
  \draw[-{Stealth[length=2mm]},thick] (0,#2)--(0,{#4+0.4}) node[below left=-1pt,font=\footnotesize]{$y$};
  \node[below left,font=\scriptsize,inner sep=1.5pt] at (0,0) {$0$};}
\newcommand{\courbet}[3]{\draw[coulCourbe,line width=1.1pt]
  plot[domain=#1:#2,samples=121] (\x,{#3});}
\newcommand{\gradt}[1]{\foreach \k in {#1}{%
  \draw (\k,0.1)--(\k,-0.1);
  \node[below,font=\scriptsize,inner sep=1pt,outer sep=2.5pt,fill=white] at (\k,0) {$\k$};}}
\newcommand{\sinusoide}[2]{%
  \begin{tikzpicture}[x=0.5cm,y=0.5cm]
    \reperet{-5}{-4}{5}{4}
    \courbet{-4}{4}{#1}
    \gradt{-4,-2,2,4}\grady{#2}
  \end{tikzpicture}}
\newcommand{\sinusoidepts}[3]{%
  \begin{tikzpicture}[x=0.5cm,y=0.5cm]
    \reperet{-5}{-4}{5}{4}
    \courbet{-4}{4}{#1}
    \gradt{-4,-2,2,4}\grady{#2}
    \foreach \u/\v in {#3}{\path[coulCourbeB] (\u,\v) node[croix]{};}
  \end{tikzpicture}}
\ShorthandsOn

\begin{remarquebox}[Comment utiliser ce cours]
Chaque partie commence par ce qu'il faut \textbf{savoir} (définitions, théorèmes), puis
vient ce qu'il faut \textbf{savoir faire} : une méthode, avec des
\textbf{exemples corrigés} faits en classe, puis des exercices
\og\textbf{À vous de jouer}\fg{} à chercher seul.

\end{remarquebox}

\partiecours{Cercle trigonométrique et radian}

Une roue, une pale d'éolienne, le bras d'un robot : tout ce qui tourne se décrit
sur un cercle de rayon $1$, avec une unité d'angle adaptée, le radian.

\begin{definitionbox}[Définition 1 --- Cercle trigonométrique]
\begin{minipage}[c]{0.66\linewidth}\raggedright
Dans un repère orthonormé $(O\,;x\,;y)$, le \textbf{cercle trigonométrique} est le
cercle de centre $O$ et de rayon $1$.

\smallskip
On le parcourt à partir du point de coordonnées $(1\,;0)$. Le \textbf{sens direct} est le sens
contraire des aiguilles d'une montre ; l'autre sens est le \textbf{sens indirect}.
\end{minipage}\hfill
\begin{minipage}[c]{0.30\linewidth}\centering
\begin{tikzpicture}[x=1.25cm,y=1.25cm]
  \cerclebase
  \arcsens{10}{80}{0.7}
  \node[coulCourbeB,font=\scriptsize] at (0.3,0.2) {direct};
\end{tikzpicture}
\end{minipage}
\end{definitionbox}

\begin{definitionbox}[Définition 2 --- Radian]
\begin{minipage}[c]{0.66\linewidth}\raggedright
Le rayon vaut $1$ : le cercle trigonométrique a pour longueur $2\pi\times 1$, soit
$2\pi$. Un tour complet mesure donc $360^{\circ}$ ou $2\pi$ \textbf{radians} (rad).

\smallskip
Sur ce cercle, la mesure en radians d'un angle est la longueur de l'arc parcouru.
\end{minipage}\hfill
\begin{minipage}[c]{0.30\linewidth}\centering
\begin{tikzpicture}[x=1.25cm,y=1.25cm]
  \cerclebase
  \draw[red,line width=1.6pt] (1,0) arc (0:60:1);
  \draw[lecture] (0,0)--({cos(60)},{sin(60)});
  \ptc{60}{$M$}
  \node[red,font=\scriptsize,right] at ({cos(30)},{sin(30)}) {arc};
\end{tikzpicture}
\end{minipage}
\end{definitionbox}

\begin{proprietebox}[Propriété 1 --- Degrés et radians]
Les mesures en degrés et en radians sont proportionnelles : $180^{\circ}$ valent $\pi$ rad.
\begin{center}
{\renewcommand{\arraystretch}{2.4}\setlength{\tabcolsep}{10pt}
\begin{tabular}{|m{1.5cm}|*{7}{>{\centering\arraybackslash}m{0.8cm}|}}
\hline
Degrés & $0$ & $30$ & $45$ & $60$ & $90$ & $180$ & $360$\\
\hline
Radians & $0$ & $\dfrac{\pi}{6}$ & $\dfrac{\pi}{4}$ & $\dfrac{\pi}{3}$ & $\dfrac{\pi}{2}$ & $\pi$ & $2\pi$\\
\hline
\end{tabular}}
\end{center}
\end{proprietebox}

\begin{methodebox}[Méthode 1 --- Convertir des degrés en radians et des radians en degrés]
\etapeexemples
Le bras d'un robot pivote. Convertir chaque rotation :
\begin{questions}[label=\alph*)]
  \item $40^{\circ}$, en radians ;
  \item $150^{\circ}$, en radians ;
  \item $\dfrac{5\pi}{12}$ rad, en degrés ;
  \item $\dfrac{7\pi}{9}$ rad, en degrés.
\end{questions}
\begin{correction}
\cas{a)}
\begin{calculs}
40^{\circ} &= 40\times\dfrac{\pi}{180}\text{ rad} &\qquad& \rem{$180^{\circ}$ valent $\pi$ rad : on multiplie par $\tfrac{\pi}{180}$}\\
40^{\circ} &= \dfrac{2\times\barre{20}\times\pi}{9\times\barre{20}}\text{ rad} && \rem{on simplifie par $20$}\\
40^{\circ} &= \dfrac{2\pi}{9}\text{ rad} &&
\end{calculs}
\cas{b)}
\begin{calculs}
150^{\circ} &= 150\times\dfrac{\pi}{180}\text{ rad} &\qquad& \rem{même calcul}\\
150^{\circ} &= \dfrac{5\times\barre{30}\times\pi}{6\times\barre{30}}\text{ rad} && \rem{on simplifie par $30$}\\
150^{\circ} &= \dfrac{5\pi}{6}\text{ rad} &&
\end{calculs}
\cas{c)}
\begin{calculs}
\dfrac{5\pi}{12}\text{ rad} &= \dfrac{5\barre{\pi}}{12}\times\dfrac{180}{\barre{\pi}}\text{ degrés} &\qquad& \rem{sens inverse : on multiplie par $\tfrac{180}{\pi}$}\\
\dfrac{5\pi}{12}\text{ rad} &= 5\times 15\text{ degrés} && \rem{$180$ divisé par $12$ vaut $15$}\\
\dfrac{5\pi}{12}\text{ rad} &= 75^{\circ} &&
\end{calculs}
\cas{d)}
\begin{calculs}
\dfrac{7\pi}{9}\text{ rad} &= \dfrac{7\barre{\pi}}{9}\times\dfrac{180}{\barre{\pi}}\text{ degrés} &\qquad& \rem{même calcul}\\
\dfrac{7\pi}{9}\text{ rad} &= 7\times 20\text{ degrés} && \rem{$180$ divisé par $9$ vaut $20$}\\
\dfrac{7\pi}{9}\text{ rad} &= 140^{\circ} &&
\end{calculs}
\end{correction}

\etape{À vous de jouer}
Une roue tourne. Convertir chaque rotation : a)~$70^{\circ}$, en radians ;
b)~$225^{\circ}$, en radians ; c)~$\dfrac{7\pi}{6}$ rad, en degrés ;
d)~$\dfrac{4\pi}{9}$ rad, en degrés.
\reponse{a)~$\dfrac{7\pi}{18}$ rad ;\; b)~$\dfrac{5\pi}{4}$ rad ;\; c)~$210^{\circ}$ ;\; d)~$80^{\circ}$.}
\end{methodebox}

\partiecours{Mesures d'un angle orienté}

\begin{definitionbox}[Définition 3 --- Mesure d'un angle orienté]
\begin{minipage}[c]{0.66\linewidth}\raggedright
On part du point de coordonnées $(1\,;0)$ et on tourne jusqu'au point $M$ du cercle
trigonométrique. La longueur parcourue est une \textbf{mesure de l'angle orienté} formé
par le demi-axe $[Ox)$ et $[OM)$ :
\begin{itemize}
  \item positive dans le sens direct ;
  \item négative dans le sens indirect.
\end{itemize}
On dit que cette mesure est le nombre réel associé au point $M$.
\end{minipage}\hfill
\begin{minipage}[c]{0.30\linewidth}\centering
\begin{tikzpicture}[x=1.25cm,y=1.25cm]
  \cerclebase
  \arcsens{0}{90}{0.72}
  \arcsens[sensB]{0}{-270}{0.45}
  \ptc{90}{$M\quad\ $}
  \node[coulCourbeB,font=\scriptsize] at (0.78,0.78) {$\frac{\pi}{2}$};
  \node[coulCourbeC,font=\scriptsize] at (-0.55,-0.55) {$-\frac{3\pi}{2}$};
\end{tikzpicture}
\end{minipage}
\end{definitionbox}

\begin{proprietebox}[Propriété 2 --- Plusieurs mesures pour un même point]
Faire un tour de plus ou de moins ramène au même point. Si $x$ est une mesure d'un angle
orienté, alors $x+2\pi$, $x-2\pi$, $x+4\pi$\ldots{} en sont aussi : toutes ses mesures
s'écrivent $x+2k\pi$, où $k$ est un entier relatif (le nombre de tours).
\end{proprietebox}

\begin{methodebox}[Méthode 2 --- Lire les mesures d'un point du cercle]
\etapeexemples
\begin{minipage}[c]{0.40\linewidth}\raggedright
Pour chaque point $A$ et $B$ de la figure, lire :
\begin{questions}[label=\alph*)]
  \item sa mesure dans le sens direct, entre $0$ et $2\pi$ ;
  \item sa mesure dans le sens indirect, entre $-2\pi$ et $0$.
\end{questions}
\end{minipage}\hfill
\begin{minipage}[c]{0.28\linewidth}\centering
Huitièmes de tour\par
\begin{tikzpicture}[x=1.25cm,y=1.25cm]
  \rayons{0,45,...,315}\cerclebase
  \ptc{135}{$A$}
\end{tikzpicture}
\end{minipage}\hfill
\begin{minipage}[c]{0.28\linewidth}\centering
Sixièmes de tour\par
\begin{tikzpicture}[x=1.25cm,y=1.25cm]
  \rayons{0,60,...,300}\cerclebase
  \ptc{300}{$B$}
\end{tikzpicture}
\end{minipage}
\begin{correction}
\cas{a)}
\explique{Pour $A$ : $3\times\dfrac{\pi}{4}$, soit $\dfrac{3\pi}{4}$.}{un huitième de tour vaut $\tfrac{\pi}{4}$ ; on en compte $3$ à partir du point $(1\,;0)$}
\explique{Pour $B$ : $5\times\dfrac{\pi}{3}$, soit $\dfrac{5\pi}{3}$.}{un sixième de tour vaut $\tfrac{\pi}{3}$ ; on en compte $5$}
\cas{b)}
\explique{Pour $A$ : $5\times\left(-\dfrac{\pi}{4}\right)$, soit $-\dfrac{5\pi}{4}$.}{dans l'autre sens, on compte $5$ huitièmes de tour}
\explique{Pour $B$ : $-\dfrac{\pi}{3}$.}{dans l'autre sens, un seul sixième de tour}
\end{correction}

\etape{À vous de jouer}
\begin{minipage}[c]{0.40\linewidth}\raggedright
Pour chaque point $C$ et $D$ de la figure, lire :
\begin{questions}[label=\alph*)]
  \item sa mesure dans le sens direct, entre $0$ et $2\pi$ ;
  \item sa mesure dans le sens indirect, entre $-2\pi$ et $0$.
\end{questions}
\end{minipage}\hfill
\begin{minipage}[c]{0.28\linewidth}\centering
Huitièmes de tour\par
\begin{tikzpicture}[x=1.25cm,y=1.25cm]
  \rayons{0,45,...,315}\cerclebase
  \ptc{225}{$C$}
\end{tikzpicture}
\end{minipage}\hfill
\begin{minipage}[c]{0.28\linewidth}\centering
Douzièmes de tour\par
\begin{tikzpicture}[x=1.25cm,y=1.25cm]
  \rayons{0,30,...,330}\cerclebase
  \ptc{150}{$D$}
\end{tikzpicture}
\end{minipage}
\reponse{a)~$C$ : $\dfrac{5\pi}{4}$ ; $D$ : $\dfrac{5\pi}{6}$ ;\; b)~$C$ : $-\dfrac{3\pi}{4}$ ; $D$ : $-\dfrac{7\pi}{6}$.}
\end{methodebox}

\begin{methodebox}[Méthode 3 --- Placer sur le cercle le point associé à un nombre réel]
\etapeexemples
Placer sur le cercle trigonométrique le point associé à chaque nombre réel :
\begin{center}
$A$ associé à $\dfrac{5\pi}{6}$ \qquad $B$ associé à $-\dfrac{2\pi}{3}$ \qquad
$C$ associé à $\dfrac{13\pi}{6}$ \qquad $D$ associé à $-\dfrac{7\pi}{2}$
\end{center}
\begin{correction}
\cas{a)}\par\nopagebreak
\begin{minipage}[c]{0.68\linewidth}\raggedright
\explique{On compte $5$ fois $\dfrac{\pi}{6}$ à partir du point $(1\,;0)$.}{positif : sens direct}
\end{minipage}\hfill
\begin{minipage}[c]{0.28\linewidth}\centering
\cerclemesure{0,30,...,330}{150}{0}{$A$}
\end{minipage}
\cas{b)}\par\nopagebreak
\begin{minipage}[c]{0.68\linewidth}\raggedright
\explique{On compte $2$ fois $\dfrac{\pi}{3}$ à partir du point $(1\,;0)$.}{négatif : sens indirect}
\end{minipage}\hfill
\begin{minipage}[c]{0.28\linewidth}\centering
\cerclemesure{0,60,...,300}{-120}{0}{$B$}
\end{minipage}
\cas{c)}\par\nopagebreak
\begin{minipage}[c]{0.68\linewidth}\raggedright
\begin{calculs}
\dfrac{13\pi}{6} &= \dfrac{12\pi}{6}+\dfrac{\pi}{6} &\quad& \rem{on fait apparaître des tours}\\
\dfrac{13\pi}{6} &= 2\pi+\dfrac{\pi}{6} && \rem{$\tfrac{12\pi}{6}$ vaut $2\pi$ : un tour}
\end{calculs}
\explique{$C$ est le point associé à $\dfrac{\pi}{6}$.}{le tour ramène au point de départ}
\end{minipage}\hfill
\begin{minipage}[c]{0.28\linewidth}\centering
\cerclemesure{0,30,...,330}{30}{0}{$C$}
\end{minipage}
\cas{d)}\par\nopagebreak
\begin{minipage}[c]{0.68\linewidth}\raggedright
\begin{calculs}
-\dfrac{7\pi}{2} &= -\dfrac{8\pi}{2}+\dfrac{\pi}{2} &\quad& \rem{$-7$ s'écrit $-8+1$}\\
-\dfrac{7\pi}{2} &= -4\pi+\dfrac{\pi}{2} && \rem{$-4\pi$ : deux tours en sens indirect}
\end{calculs}
\explique{$D$ est le point associé à $\dfrac{\pi}{2}$.}{les deux tours ramènent au point de départ}
\end{minipage}\hfill
\begin{minipage}[c]{0.28\linewidth}\centering
\cerclemesure{0,90,180,270}{90}{0}{$D\quad\ $}
\end{minipage}
\end{correction}

\etape{À vous de jouer}
Placer sur le cercle trigonométrique le point associé à chaque nombre réel :
\begin{center}
$E$ associé à $\dfrac{2\pi}{3}$ \qquad $F$ associé à $-\dfrac{3\pi}{4}$ \qquad
$G$ associé à $\dfrac{17\pi}{6}$ \qquad $H$ associé à $-\dfrac{5\pi}{2}$
\end{center}
\reponse{$G$ est le point associé à $\dfrac{5\pi}{6}$ ;\; $H$ est le point associé à $-\dfrac{\pi}{2}$.}
\end{methodebox}

\partiecours{Mesure principale d'un angle orienté}

Un même point a une infinité de mesures : on en choisit une, toujours la même.

\begin{definitionbox}[Définition 4 --- Mesure principale]
\begin{minipage}[c]{0.66\linewidth}\raggedright
La \textbf{mesure principale} d'un angle orienté est sa seule mesure comprise dans
l'intervalle $\intof{-\pi}{\pi}$.

\smallskip
Les points du haut du cercle ont une mesure principale positive, ceux du bas une mesure
principale négative.
\end{minipage}\hfill
\begin{minipage}[c]{0.30\linewidth}\centering
\begin{tikzpicture}[x=1.25cm,y=1.25cm]
  \cerclebase
  \draw[coulCourbeB,line width=1.5pt] (1,0) arc (0:180:1);
  \draw[coulCourbeC,line width=1.5pt] (1,0) arc (0:-180:1);
  \node[coulCourbeB,font=\scriptsize] at (-1.2,0.18) {$\pi$};
  \node[coulCourbeB,font=\scriptsize] at (0.45,0.45) {$+$};
  \node[coulCourbeC,font=\scriptsize] at (0.45,-0.45) {$-$};
\end{tikzpicture}
\end{minipage}
\end{definitionbox}

\begin{methodebox}[Méthode 4 --- Déterminer la mesure principale d'un angle orienté]
\etapeexemples
Déterminer la mesure principale de chaque angle orienté, de mesure :
\begin{center}
a)~$\dfrac{8\pi}{3}$ \qquad b)~$-\dfrac{7\pi}{6}$ \qquad c)~$\dfrac{29\pi}{4}$ \qquad d)~$7\pi$
\end{center}
\begin{correction}
\cas{a)}
\begin{calculs}
2\pi &= \dfrac{6\pi}{3} &\qquad& \rem{un tour, écrit avec le dénominateur $3$}\\
\dfrac{8\pi}{3}-\dfrac{6\pi}{3} &= \dfrac{2\pi}{3} && \rem{on retire un tour}
\end{calculs}
\explique{La mesure principale est $\dfrac{2\pi}{3}$.}{$\tfrac{2\pi}{3}$ est bien dans $\intof{-\pi}{\pi}$}
\cas{b)}
\begin{calculs}
2\pi &= \dfrac{12\pi}{6} &\qquad& \rem{un tour, écrit avec le dénominateur $6$}\\
-\dfrac{7\pi}{6}+\dfrac{12\pi}{6} &= \dfrac{5\pi}{6} && \rem{mesure négative : on ajoute un tour}
\end{calculs}
\explique{La mesure principale est $\dfrac{5\pi}{6}$.}{$\tfrac{5\pi}{6}$ est bien dans $\intof{-\pi}{\pi}$}
\cas{c)}
\begin{calculs}
2\pi &= \dfrac{8\pi}{4} &\qquad& \rem{un tour, écrit avec le dénominateur $4$}\\
4\times 2\pi &= \dfrac{32\pi}{4} && \rem{$32$ : le multiple de $8$ le plus proche de $29$}\\
\dfrac{29\pi}{4}-\dfrac{32\pi}{4} &= -\dfrac{3\pi}{4} && \rem{on retire $4$ tours}
\end{calculs}
\explique{La mesure principale est $-\dfrac{3\pi}{4}$.}{$-\tfrac{3\pi}{4}$ est bien dans $\intof{-\pi}{\pi}$ ; avec $24$, on trouverait $\tfrac{5\pi}{4}$, trop grand}
\cas{d)}
\begin{calculs}
7\pi-6\pi &= \pi &\qquad& \rem{$6\pi$ : $3$ tours}
\end{calculs}
\explique{La mesure principale est $\pi$.}{$\pi$ est dans $\intof{-\pi}{\pi}$, mais pas $-\pi$}
\end{correction}

\etape{À vous de jouer}
Déterminer la mesure principale de chaque angle orienté, de mesure :
\begin{center}
a)~$\dfrac{11\pi}{3}$ \qquad b)~$-\dfrac{13\pi}{6}$ \qquad c)~$\dfrac{35\pi}{4}$ \qquad d)~$9\pi$
\end{center}
\reponse{a)~$-\dfrac{\pi}{3}$ ;\; b)~$-\dfrac{\pi}{6}$ ;\; c)~$\dfrac{3\pi}{4}$ ;\; d)~$\pi$.}
\end{methodebox}

\partiecours{Cosinus et sinus d'un nombre réel}

\begin{definitionbox}[Définition 5 --- Cosinus et sinus d'un nombre réel]
\begin{minipage}[c]{0.66\linewidth}\raggedright
Soit $M$ le point du cercle trigonométrique associé au nombre réel $x$.
\begin{itemize}
  \item Le \textbf{cosinus} de $x$, noté $\cos(x)$, est l'abscisse de $M$.
  \item Le \textbf{sinus} de $x$, noté $\sin(x)$, est l'ordonnée de $M$.
\end{itemize}
Entre $0$ et $\dfrac{\pi}{2}$, on retrouve le cosinus et le sinus du triangle rectangle :
dans le triangle $OHM$, l'hypoténuse $OM$ vaut $1$.
\end{minipage}\hfill
\begin{minipage}[c]{0.30\linewidth}\centering
\begin{tikzpicture}[x=1.7cm,y=1.7cm]
  \cerclebase
  \draw[black!70] (0,0)--({cos(50)},{sin(50)});
  \draw[lecture] ({cos(50)},{sin(50)})--({cos(50)},0);
  \draw[lecture] ({cos(50)},{sin(50)})--(0,{sin(50)});
  \draw[coulCourbeB,line width=1.6pt] (0,0)--({cos(50)},0);
  \draw[coulCourbeC,line width=1.6pt] (0,0)--(0,{sin(50)});
  \ptc{50}{$M$}
  \node[font=\scriptsize,above right,inner sep=1pt] at ({cos(50)},0) {$H$};
  \node[coulCourbeB,font=\scriptsize,below] at ({cos(50)/2},0) {$\cos(x)$};
  \node[coulCourbeC,font=\scriptsize,left] at (0,{sin(50)/2}) {$\sin(x)$};
\end{tikzpicture}
\end{minipage}
\end{definitionbox}

\begin{proprietebox}[Propriété 3 --- Encadrement et relation]
Pour tout nombre réel $x$ :
\begin{center}
$-1\le\cos(x)\le 1$ \qquad $-1\le\sin(x)\le 1$ \qquad $\cos^{2}(x)+\sin^{2}(x)=1$
\end{center}
où $\cos^{2}(x)$ désigne $\big(\cos(x)\big)^{2}$. Et comme un tour ramène au même point :
\begin{center}
$\cos(x+2\pi)=\cos(x)$ \qquad $\sin(x+2\pi)=\sin(x)$
\end{center}
\demo{Le point $M$ est sur un cercle de rayon $1$ : ses coordonnées sont entre $-1$ et $1$.
Le triangle $OHM$ est rectangle en $H$ :}
\begin{calculs}
OH^{2}+HM^{2} &= OM^{2} &\qquad& \rem{théorème de Pythagore}\\
\cos^{2}(x)+\sin^{2}(x) &= 1 && \rem{$OM$ vaut $1$}
\end{calculs}
\end{proprietebox}

\begin{proprietebox}[Propriété 4 --- Valeurs remarquables]
\begin{minipage}[c]{0.60\linewidth}\centering
{\renewcommand{\arraystretch}{2.4}\setlength{\tabcolsep}{7pt}
\begin{tabular}{|>{\centering\arraybackslash}m{1.2cm}|*{6}{>{\centering\arraybackslash}m{0.65cm}|}}
\hline
$x$ & $0$ & $\dfrac{\pi}{6}$ & $\dfrac{\pi}{4}$ & $\dfrac{\pi}{3}$ & $\dfrac{\pi}{2}$ & $\pi$\\
\hline
$\cos(x)$ & $1$ & $\dfrac{\sqrt{3}}{2}$ & $\dfrac{\sqrt{2}}{2}$ & $\dfrac{1}{2}$ & $0$ & $-1$\\
\hline
$\sin(x)$ & $0$ & $\dfrac{1}{2}$ & $\dfrac{\sqrt{2}}{2}$ & $\dfrac{\sqrt{3}}{2}$ & $1$ & $0$\\
\hline
\end{tabular}}

\medskip
{\raggedright À retenir avec l'image du quart de cercle : le cosinus diminue de $1$ à $0$
quand $x$ augmente, le sinus augmente de $0$ à $1$.\par}
\end{minipage}\hfill
\begin{minipage}[c]{0.36\linewidth}\centering
\begin{tikzpicture}[x=4cm,y=4cm]
  \draw[-{Stealth[length=2mm]},thick] (0,0)--(1.2,0);
  \draw[-{Stealth[length=2mm]},thick] (0,0)--(0,1.2);
  \draw[line width=0.9pt] (1,0) arc (0:90:1);
  \foreach \a/\l/\c/\s in {30/\frac{\pi}{6}/\frac{\sqrt{3}}{2}/\frac{1}{2},
      45/\frac{\pi}{4}/\frac{\sqrt{2}}{2}/\frac{\sqrt{2}}{2},
      60/\frac{\pi}{3}/\frac{1}{2}/\frac{\sqrt{3}}{2}}{%
    \draw[lecture] ({cos(\a)},0)--({cos(\a)},{sin(\a)})--(0,{sin(\a)});
    \draw[grille] (0,0)--({cos(\a)},{sin(\a)});
    \path[coulCourbeB] ({cos(\a)},{sin(\a)}) node[croix]{};
    \node[font=\scriptsize] at ({1.12*cos(\a)},{1.12*sin(\a)}) {$\l$};
    \node[coulCourbeB,font=\scriptsize,below,inner sep=2pt] at ({cos(\a)},0) {$\c$};
    \node[coulCourbeC,font=\scriptsize,left,inner sep=2pt] at (0,{sin(\a)}) {$\s$};}
  \node[font=\scriptsize,below,inner sep=2pt] at (1,0) {$1$};
  \node[font=\scriptsize,left,inner sep=2pt] at (0,1) {$1$};
  \node[font=\scriptsize,below left,inner sep=2pt] at (0,0) {$O$};
\end{tikzpicture}
\end{minipage}
\end{proprietebox}

\begin{methodebox}[Méthode 5 --- Déterminer un cosinus ou un sinus à l'aide de la mesure principale]
\etapeexemples
Déterminer la valeur exacte de :
\begin{center}
a)~$\cos\left(\dfrac{\pi}{4}\right)$ \qquad b)~$\sin\left(\dfrac{13\pi}{6}\right)$ \qquad
c)~$\cos(9\pi)$ \qquad d)~$\sin\left(-\dfrac{3\pi}{2}\right)$
\end{center}
\begin{correction}
\cas{a)}
\explique{$\cos\left(\dfrac{\pi}{4}\right)=\dfrac{\sqrt{2}}{2}$.}{lecture du tableau}
\cas{b)}
\begin{calculs}
\dfrac{13\pi}{6}-\dfrac{12\pi}{6} &= \dfrac{\pi}{6} &\qquad& \rem{on retire un tour : même point}\\
\sin\left(\dfrac{13\pi}{6}\right) &= \sin\left(\dfrac{\pi}{6}\right) && \rem{même point, même sinus}\\
\sin\left(\dfrac{13\pi}{6}\right) &= \dfrac{1}{2} && \rem{lecture du tableau}
\end{calculs}
\cas{c)}
\begin{calculs}
9\pi-8\pi &= \pi &\qquad& \rem{on retire $4$ tours}\\
\cos(9\pi) &= \cos(\pi) && \rem{même point, même cosinus}\\
\cos(9\pi) &= -1 && \rem{lecture du tableau}
\end{calculs}
\cas{d)}
\begin{calculs}
-\dfrac{3\pi}{2}+\dfrac{4\pi}{2} &= \dfrac{\pi}{2} &\qquad& \rem{on ajoute un tour}\\
\sin\left(-\dfrac{3\pi}{2}\right) &= \sin\left(\dfrac{\pi}{2}\right) && \rem{même point, même sinus}\\
\sin\left(-\dfrac{3\pi}{2}\right) &= 1 && \rem{lecture du tableau}
\end{calculs}
\end{correction}

\etape{À vous de jouer}
Déterminer la valeur exacte de :
\begin{center}
a)~$\cos\left(\dfrac{\pi}{3}\right)$ \qquad b)~$\sin\left(\dfrac{25\pi}{6}\right)$ \qquad
c)~$\cos(7\pi)$ \qquad d)~$\sin\left(-\dfrac{7\pi}{2}\right)$
\end{center}
\reponse{a)~$\dfrac{1}{2}$ ;\; b)~$\dfrac{1}{2}$ ;\; c)~$-1$ ;\; d)~$1$.}
\end{methodebox}

\begin{methodebox}[Méthode 6 --- Calculer une expression avec des cosinus et des sinus]
\etapeexemples
Calculer la valeur exacte de :
\begin{questions}[label=\alph*)]
  \item $f(x)=4\cos(x)-2\sin(x)$, pour $x=\dfrac{\pi}{6}$ ;
  \item $g(x)=\sin^{2}(x)$, pour $x=\dfrac{\pi}{4}$ ;
  \item $h(x)=\cos(2x)$, pour $x=\dfrac{\pi}{6}$.
\end{questions}
\begin{correction}
\cas{a)}
\begin{calculs}
f\left(\dfrac{\pi}{6}\right) &= 4\cos\left(\dfrac{\pi}{6}\right)-2\sin\left(\dfrac{\pi}{6}\right) &\qquad& \rem{on remplace $x$ par $\tfrac{\pi}{6}$}\\
f\left(\dfrac{\pi}{6}\right) &= 4\times\dfrac{\sqrt{3}}{2}-2\times\dfrac{1}{2} && \rem{valeurs du tableau}\\
f\left(\dfrac{\pi}{6}\right) &= 2\sqrt{3}-1 && \rem{$4$ divisé par $2$ vaut $2$}
\end{calculs}
\cas{b)}
\begin{calculs}
g\left(\dfrac{\pi}{4}\right) &= \left(\sin\left(\dfrac{\pi}{4}\right)\right)^{2} &\qquad& \rem{$\sin^{2}$ : le sinus, au carré}\\
g\left(\dfrac{\pi}{4}\right) &= \left(\dfrac{\sqrt{2}}{2}\right)^{2} && \rem{valeur du tableau}\\
g\left(\dfrac{\pi}{4}\right) &= \dfrac{2}{4} && \rem{$\big(\sqrt{2}\big)^{2}$ vaut $2$ et $2^{2}$ vaut $4$}\\
g\left(\dfrac{\pi}{4}\right) &= \dfrac{1}{2} &&
\end{calculs}
\cas{c)}
\begin{calculs}
h\left(\dfrac{\pi}{6}\right) &= \cos\left(2\times\dfrac{\pi}{6}\right) &\qquad& \rem{on multiplie d'abord, on prend le cosinus ensuite}\\
h\left(\dfrac{\pi}{6}\right) &= \cos\left(\dfrac{\pi}{3}\right) && \rem{$2\times\tfrac{\pi}{6}$ vaut $\tfrac{\pi}{3}$}\\
h\left(\dfrac{\pi}{6}\right) &= \dfrac{1}{2} && \rem{valeur du tableau}
\end{calculs}
\end{correction}

\etape{À vous de jouer}
Calculer la valeur exacte de :
\begin{questions}[label=\alph*)]
  \item $k(x)=2\cos(x)+6\sin(x)$, pour $x=\dfrac{\pi}{6}$ ;
  \item $m(x)=\cos^{2}(x)$, pour $x=\dfrac{\pi}{6}$ ;
  \item $n(x)=\sin(2x)$, pour $x=\dfrac{\pi}{4}$.
\end{questions}
\reponse{a)~$\sqrt{3}+3$ ;\; b)~$\dfrac{3}{4}$ ;\; c)~$1$.}
\end{methodebox}

\partiecours{Angles associés}

Le tableau s'arrête à $\pi$ et ne donne rien pour $\dfrac{2\pi}{3}$ ou $-\dfrac{\pi}{4}$ :
ces points sont les symétriques de points connus, leurs coordonnées s'en déduisent.

\begin{proprietebox}[Propriété 5 --- Angles associés]
Pour tout nombre réel $x$, on lit sur le cercle :

\medskip
\begin{minipage}[t]{0.32\linewidth}\centering
\figassociee{-35}{$-x$}{\draw[coulCourbeC,line width=0.8pt] (-1.2,0)--(1.2,0);\lien{35}{-35}\liery{35}\liery{-35}}\par
{\small symétrie d'axe $(Ox)$ :\\ même cosinus, sinus opposés}\par\smallskip
$\cos(-x)=\cos(x)$\par
$\sin(-x)=-\sin(x)$
\end{minipage}\hfill
\begin{minipage}[t]{0.32\linewidth}\centering
\figassociee{145}{$\pi-x$}{\draw[coulCourbeC,line width=0.8pt] (0,-1.2)--(0,1.2);\lien{35}{145}\lierx{35}\lierx{145}}\par
{\small symétrie d'axe $(Oy)$ :\\ même sinus, cosinus opposés}\par\smallskip
$\cos(\pi-x)=-\cos(x)$\par
$\sin(\pi-x)=\sin(x)$
\end{minipage}\hfill
\begin{minipage}[t]{0.32\linewidth}\centering
\figassociee{215}{$\pi+x$}{\lien{35}{215}\lierx{35}\liery{35}\lierx{215}\liery{215}}\par
{\small symétrie de centre $O$ :\\ cosinus opposés, sinus opposés}\par\smallskip
$\cos(\pi+x)=-\cos(x)$\par
$\sin(\pi+x)=-\sin(x)$
\end{minipage}

\medskip
\bigskip
\hfill
\begin{minipage}[t]{0.32\linewidth}\centering
\figassociee{55}{$\frac{\pi}{2}-x\quad$}{\draw[coulCourbeC,line width=0.8pt] (-1.1,-1.1)--(1.1,1.1);\lien{35}{55}\liery{35}\lierx{55}}\par
{\small symétrie d'axe la droite $y=x$ :\\ cosinus et sinus s'échangent}\par\smallskip
$\cos\left(\dfrac{\pi}{2}-x\right)=\sin(x)$\par\smallskip
$\sin\left(\dfrac{\pi}{2}-x\right)=\cos(x)$
\end{minipage}\hspace{1.5cm}
\begin{minipage}[t]{0.32\linewidth}\centering
\figassociee{125}{$\frac{\pi}{2}+x$}{\draw[lecture] (0,0)--({cos(35)},{sin(35)}); \draw[lecture] (0,0)--({cos(125)},{sin(125)});\liery{35}\lierx{125}}\par
{\small quart de tour de centre $O$ :\\ ils s'échangent, un signe change}\par\smallskip
$\cos\left(\dfrac{\pi}{2}+x\right)=-\sin(x)$\par\smallskip
$\sin\left(\dfrac{\pi}{2}+x\right)=\cos(x)$
\end{minipage}
\hfill\null

\medskip
Ces formules ne s'apprennent pas par cœur une à une : on dessine la figure et on regarde
si l'abscisse et l'ordonnée changent de signe ou s'échangent.
\end{proprietebox}

\begin{methodebox}[Méthode 7 --- Calculer un cosinus ou un sinus à l'aide d'un angle associé]
\etapeexemples
Déterminer la valeur exacte de :
\begin{center}
a)~$\cos\left(\dfrac{3\pi}{4}\right)$ \qquad b)~$\sin\left(\dfrac{2\pi}{3}\right)$ \qquad
c)~$\sin\left(\dfrac{7\pi}{6}\right)$ \qquad d)~$\cos\left(-\dfrac{\pi}{6}\right)$
\end{center}
\begin{correction}
\cas{a)}
\begin{calculs}
\cos\left(\dfrac{3\pi}{4}\right) &= \cos\left(\pi-\dfrac{\pi}{4}\right) &\qquad& \rem{$\tfrac{4\pi}{4}-\tfrac{\pi}{4}$ vaut $\tfrac{3\pi}{4}$}\\
\cos\left(\dfrac{3\pi}{4}\right) &= -\cos\left(\dfrac{\pi}{4}\right) && \rem{$\cos(\pi-x)=-\cos(x)$}\\
\cos\left(\dfrac{3\pi}{4}\right) &= -\dfrac{\sqrt{2}}{2} && \rem{valeur du tableau}
\end{calculs}
\cas{b)}
\begin{calculs}
\sin\left(\dfrac{2\pi}{3}\right) &= \sin\left(\pi-\dfrac{\pi}{3}\right) &\qquad& \rem{$\tfrac{3\pi}{3}-\tfrac{\pi}{3}$ vaut $\tfrac{2\pi}{3}$}\\
\sin\left(\dfrac{2\pi}{3}\right) &= \sin\left(\dfrac{\pi}{3}\right) && \rem{$\sin(\pi-x)=\sin(x)$}\\
\sin\left(\dfrac{2\pi}{3}\right) &= \dfrac{\sqrt{3}}{2} && \rem{valeur du tableau}
\end{calculs}
\cas{c)}
\begin{calculs}
\sin\left(\dfrac{7\pi}{6}\right) &= \sin\left(\pi+\dfrac{\pi}{6}\right) &\qquad& \rem{$\tfrac{6\pi}{6}+\tfrac{\pi}{6}$ vaut $\tfrac{7\pi}{6}$}\\
\sin\left(\dfrac{7\pi}{6}\right) &= -\sin\left(\dfrac{\pi}{6}\right) && \rem{$\sin(\pi+x)=-\sin(x)$}\\
\sin\left(\dfrac{7\pi}{6}\right) &= -\dfrac{1}{2} && \rem{valeur du tableau}
\end{calculs}
\cas{d)}
\begin{calculs}
\cos\left(-\dfrac{\pi}{6}\right) &= \cos\left(\dfrac{\pi}{6}\right) &\qquad& \rem{$\cos(-x)=\cos(x)$}\\
\cos\left(-\dfrac{\pi}{6}\right) &= \dfrac{\sqrt{3}}{2} && \rem{valeur du tableau}
\end{calculs}
\end{correction}

\etape{À vous de jouer}
Déterminer la valeur exacte de :
\begin{center}
a)~$\cos\left(\dfrac{5\pi}{6}\right)$ \qquad b)~$\sin\left(\dfrac{3\pi}{4}\right)$ \qquad
c)~$\sin\left(\dfrac{5\pi}{4}\right)$ \qquad d)~$\cos\left(-\dfrac{\pi}{4}\right)$
\end{center}
\reponse{a)~$-\dfrac{\sqrt{3}}{2}$ ;\; b)~$\dfrac{\sqrt{2}}{2}$ ;\; c)~$-\dfrac{\sqrt{2}}{2}$ ;\; d)~$\dfrac{\sqrt{2}}{2}$.}
\end{methodebox}

\begin{methodebox}[Méthode 8 --- Simplifier une expression à l'aide des angles associés]
\etapeexemples
Pour tout nombre réel $x$, simplifier :
\begin{center}
$A(x)=\cos(\pi-x)+\cos(-x)$ \qquad
$B(x)=\sin\left(\dfrac{\pi}{2}-x\right)+\cos(\pi+x)+\sin(-x)$\\[4pt]
$C(x)=\cos\left(\dfrac{\pi}{2}+x\right)+3\sin(\pi-x)$
\end{center}
\begin{correction}
\cas{Expression $A$}
\begin{calculs}
A(x) &= \cos(\pi-x)+\cos(-x) &\qquad& \rem{on recopie}\\
A(x) &= -\cos(x)+\cos(x) && \rem{on remplace chaque terme, d'après la propriété 5}\\
A(x) &= 0 &&
\end{calculs}
\cas{Expression $B$}
\begin{calculs}
B(x) &= \sin\left(\dfrac{\pi}{2}-x\right)+\cos(\pi+x)+\sin(-x) &\qquad& \rem{on recopie}\\
B(x) &= \cos(x)-\cos(x)-\sin(x) && \rem{on remplace chaque terme}\\
B(x) &= -\sin(x) && \rem{les cosinus s'annulent}
\end{calculs}
\cas{Expression $C$}
\begin{calculs}
C(x) &= \cos\left(\dfrac{\pi}{2}+x\right)+3\sin(\pi-x) &\qquad& \rem{on recopie}\\
C(x) &= -\sin(x)+3\sin(x) && \rem{on remplace chaque terme}\\
C(x) &= 2\sin(x) && \rem{$-1+3$ vaut $2$}
\end{calculs}
\end{correction}

\etape{À vous de jouer}
Pour tout nombre réel $x$, simplifier :
\begin{center}
$D(x)=\sin(\pi+x)+\sin(\pi-x)$ \qquad
$E(x)=\cos\left(\dfrac{\pi}{2}-x\right)+\cos(\pi+x)+\cos(-x)$\\[4pt]
$F(x)=\sin\left(\dfrac{\pi}{2}+x\right)+4\cos(\pi-x)$
\end{center}
\reponse{$D(x)=0$ ;\; $E(x)=\sin(x)$ ;\; $F(x)=-3\cos(x)$.}
\end{methodebox}

\begin{proprietebox}[Le cercle complet et son tableau de valeurs]
Les symétries donnent seize points, chacun avec sa mesure principale. Les pointillés
verticaux relient les points de même cosinus, les pointillés horizontaux ceux de même sinus.
Ce cercle est l'outil de tout le chapitre : \textbf{il faut savoir le refaire}.
\begin{center}
\begin{tikzpicture}[x=3.7cm,y=3.7cm]
  \foreach \a in {30,45,60}{%
    \draw[lecture] ({cos(\a)},{-sin(\a)})--({cos(\a)},{sin(\a)});
    \draw[lecture] ({-cos(\a)},{-sin(\a)})--({-cos(\a)},{sin(\a)});
    \draw[lecture] ({-cos(\a)},{sin(\a)})--({cos(\a)},{sin(\a)});
    \draw[lecture] ({-cos(\a)},{-sin(\a)})--({cos(\a)},{-sin(\a)});}
  \cerclebase
  \path[coulCourbeB] (1,0) node[croix]{} (0,1) node[croix]{} (-1,0) node[croix]{} (0,-1) node[croix]{};
  \node[font=\footnotesize,inner sep=1pt] at (1.1,0.09) {$0$};
  \node[font=\footnotesize,inner sep=1pt] at (-1.1,0.09) {$\pi$};
  \node[font=\footnotesize,inner sep=1pt] at (-0.12,1.12) {$\frac{\pi}{2}$};
  \node[font=\footnotesize,inner sep=1pt] at (-0.16,-1.12) {$-\frac{\pi}{2}$};
  \foreach \a/\l in {30/\frac{\pi}{6},45/\frac{\pi}{4},60/\frac{\pi}{3},
     120/\frac{2\pi}{3},135/\frac{3\pi}{4},150/\frac{5\pi}{6},
     -30/-\frac{\pi}{6},-45/-\frac{\pi}{4},-60/-\frac{\pi}{3},
     -120/-\frac{2\pi}{3},-135/-\frac{3\pi}{4},-150/-\frac{5\pi}{6}}{%
    \path[coulCourbeB] ({cos(\a)},{sin(\a)}) node[croix]{};
    \node[font=\footnotesize,inner sep=1pt] at ({1.13*cos(\a)},{1.13*sin(\a)}) {$\l$};}
  \foreach \v/\l in {0.5/\frac{1}{2},0.866/\frac{\sqrt{3}}{2},-0.5/-\frac{1}{2},-0.866/-\frac{\sqrt{3}}{2}}{%
    \draw (\v,0.02)--(\v,-0.02);
    \node[below,font=\footnotesize,inner sep=1pt,outer sep=2pt,coulCourbeB,fill=white] at (\v,0) {$\l$};
    \draw (0.02,\v)--(-0.02,\v);
    \node[left,font=\footnotesize,inner sep=1pt,outer sep=2pt,coulCourbeC,fill=white] at (0,\v) {$\l$};}
  \foreach \v/\l in {0.707/\frac{\sqrt{2}}{2},-0.707/-\frac{\sqrt{2}}{2}}{%
    \draw (\v,0.02)--(\v,-0.02);
    \node[above,font=\footnotesize,inner sep=1pt,outer sep=2pt,coulCourbeB,fill=white] at (\v,0) {$\l$};
    \draw (0.02,\v)--(-0.02,\v);
    \node[right,font=\footnotesize,inner sep=1pt,outer sep=2pt,coulCourbeC,fill=white] at (0,\v) {$\l$};}
\end{tikzpicture}
\end{center}
\begin{center}
{\renewcommand{\arraystretch}{2.4}\setlength{\tabcolsep}{4pt}
\begin{tabular}{|>{\centering\arraybackslash}m{1.2cm}|*{9}{>{\centering\arraybackslash}m{1.15cm}|}}
\hline
$x$ & $0$ & $\dfrac{\pi}{6}$ & $\dfrac{\pi}{4}$ & $\dfrac{\pi}{3}$ & $\dfrac{\pi}{2}$ & $\dfrac{2\pi}{3}$ & $\dfrac{3\pi}{4}$ & $\dfrac{5\pi}{6}$ & $\pi$\\
\hline
$\cos(x)$ & $1$ & $\dfrac{\sqrt{3}}{2}$ & $\dfrac{\sqrt{2}}{2}$ & $\dfrac{1}{2}$ & $0$ & $-\dfrac{1}{2}$ & $-\dfrac{\sqrt{2}}{2}$ & $-\dfrac{\sqrt{3}}{2}$ & $-1$\\
\hline
$\sin(x)$ & $0$ & $\dfrac{1}{2}$ & $\dfrac{\sqrt{2}}{2}$ & $\dfrac{\sqrt{3}}{2}$ & $1$ & $\dfrac{\sqrt{3}}{2}$ & $\dfrac{\sqrt{2}}{2}$ & $\dfrac{1}{2}$ & $0$\\
\hline
\end{tabular}

\medskip
\renewcommand{\arraystretch}{2.4}
\begin{tabular}{|>{\centering\arraybackslash}m{1.2cm}|*{7}{>{\centering\arraybackslash}m{1.15cm}|}}
\hline
$x$ & $-\dfrac{\pi}{6}$ & $-\dfrac{\pi}{4}$ & $-\dfrac{\pi}{3}$ & $-\dfrac{\pi}{2}$ & $-\dfrac{2\pi}{3}$ & $-\dfrac{3\pi}{4}$ & $-\dfrac{5\pi}{6}$\\
\hline
$\cos(x)$ & $\dfrac{\sqrt{3}}{2}$ & $\dfrac{\sqrt{2}}{2}$ & $\dfrac{1}{2}$ & $0$ & $-\dfrac{1}{2}$ & $-\dfrac{\sqrt{2}}{2}$ & $-\dfrac{\sqrt{3}}{2}$\\
\hline
$\sin(x)$ & $-\dfrac{1}{2}$ & $-\dfrac{\sqrt{2}}{2}$ & $-\dfrac{\sqrt{3}}{2}$ & $-1$ & $-\dfrac{\sqrt{3}}{2}$ & $-\dfrac{\sqrt{2}}{2}$ & $-\dfrac{1}{2}$\\
\hline
\end{tabular}}
\end{center}
\end{proprietebox}

\partiecours{Équations avec des cosinus et des sinus}

Résoudre $\cos(x)=a$, c'est chercher les points du cercle d'abscisse $a$, sur une droite
verticale ; résoudre $\sin(x)=a$, c'est chercher les points d'ordonnée $a$, sur une droite
horizontale. Ces équations se résolvent avec le \textbf{cercle complet} de la partie V sous
les yeux : il est donné dans les exercices ; en DS, il est à compléter.

\begin{methodebox}[Méthode 9 --- Résoudre $\cos(x)=a$ ou $\sin(x)=a$ sur $\intof{-\pi}{\pi}$]
\etapeexemples
Résoudre sur l'intervalle $\intof{-\pi}{\pi}$, à l'aide du cercle complet :
\begin{center}
a)~$\cos(x)=\dfrac{1}{2}$ \qquad b)~$\sin(x)=\dfrac{\sqrt{2}}{2}$ \qquad
c)~$\cos(x)=-\dfrac{\sqrt{3}}{2}$\\[4pt]
d)~$\sin(x)=-\dfrac{1}{2}$ \qquad e)~$\cos(x)=-1$ \qquad f)~$\sin(x)=2$
\end{center}
\begin{correction}
\cas{a)}\par\nopagebreak
\begin{minipage}[c]{0.68\linewidth}\raggedright
\explique{On place $\dfrac{1}{2}$ sur l'axe des abscisses et on trace la droite verticale.}{le cosinus se lit sur l'axe des abscisses}
\explique{Elle coupe le cercle en deux points, associés à $\dfrac{\pi}{3}$ et $-\dfrac{\pi}{3}$.}{lecture sur le cercle complet}
\explique{$S=\ens{-\dfrac{\pi}{3}\,;\dfrac{\pi}{3}}$.}{une solution par point}
\end{minipage}\hfill
\begin{minipage}[c]{0.28\linewidth}\centering
\equacos{60}{$\frac{\pi}{3}$}{$-\frac{\pi}{3}$}{$\frac{1}{2}$}
\end{minipage}
\cas{b)}\par\nopagebreak
\begin{minipage}[c]{0.68\linewidth}\raggedright
\explique{On place $\dfrac{\sqrt{2}}{2}$ sur l'axe des ordonnées et on trace la droite horizontale.}{le sinus se lit sur l'axe des ordonnées}
\explique{Elle coupe le cercle en deux points, associés à $\dfrac{\pi}{4}$ et $\dfrac{3\pi}{4}$.}{lecture sur le cercle complet}
\explique{$S=\ens{\dfrac{\pi}{4}\,;\dfrac{3\pi}{4}}$.}{une solution par point}
\end{minipage}\hfill
\begin{minipage}[c]{0.28\linewidth}\centering
\equasin{45}{135}{$\frac{\pi}{4}$}{$\frac{3\pi}{4}$}{$\frac{\sqrt{2}}{2}$}
\end{minipage}
\cas{c)}\par\nopagebreak
\begin{minipage}[c]{0.68\linewidth}\raggedright
\explique{On place $-\dfrac{\sqrt{3}}{2}$ sur l'axe des abscisses et on trace la droite verticale.}{valeur négative : à gauche de $O$}
\explique{Elle coupe le cercle en deux points, associés à $\dfrac{5\pi}{6}$ et $-\dfrac{5\pi}{6}$.}{on les lit directement sur le cercle complet}
\explique{$S=\ens{-\dfrac{5\pi}{6}\,;\dfrac{5\pi}{6}}$.}{une solution par point}
\end{minipage}\hfill
\begin{minipage}[c]{0.28\linewidth}\centering
\equacos{150}{$\frac{5\pi}{6}$}{$-\frac{5\pi}{6}$}{$-\frac{\sqrt{3}}{2}$}
\end{minipage}
\cas{d)}\par\nopagebreak
\begin{minipage}[c]{0.68\linewidth}\raggedright
\explique{On place $-\dfrac{1}{2}$ sur l'axe des ordonnées et on trace la droite horizontale.}{valeur négative : au-dessous de $O$}
\explique{Elle coupe le cercle en deux points, associés à $-\dfrac{\pi}{6}$ et $-\dfrac{5\pi}{6}$.}{on les lit directement sur le cercle complet}
\explique{$S=\ens{-\dfrac{5\pi}{6}\,;-\dfrac{\pi}{6}}$.}{une solution par point}
\end{minipage}\hfill
\begin{minipage}[c]{0.28\linewidth}\centering
\equasin{-30}{-150}{$-\frac{\pi}{6}$}{$-\frac{5\pi}{6}$}{$-\frac{1}{2}$}
\end{minipage}
\cas{e)}\par\nopagebreak
\begin{minipage}[c]{0.68\linewidth}\raggedright
\explique{La droite verticale passant par $-1$ touche le cercle en un seul point.}{le point le plus à gauche du cercle}
\explique{$S=\ens{\pi}$.}{une seule solution}
\end{minipage}\hfill
\begin{minipage}[c]{0.28\linewidth}\centering
\cerclept{180}{$\pi$}{\verticale{180}}
\end{minipage}
\cas{f)}
\explique{Il n'y a aucune solution.}{la droite horizontale passant par $2$ ne rencontre pas le cercle : un sinus ne dépasse jamais $1$}
\end{correction}

\etape{À vous de jouer}
Résoudre sur l'intervalle $\intof{-\pi}{\pi}$, à l'aide du cercle complet :
\begin{center}
a)~$\cos(x)=\dfrac{\sqrt{2}}{2}$ \qquad b)~$\sin(x)=\dfrac{\sqrt{3}}{2}$ \qquad
c)~$\cos(x)=-\dfrac{1}{2}$\\[4pt]
d)~$\sin(x)=-\dfrac{\sqrt{2}}{2}$ \qquad e)~$\cos(x)=1$ \qquad f)~$\cos(x)=-3$
\end{center}
\reponse{a)~$\ens{-\dfrac{\pi}{4}\,;\dfrac{\pi}{4}}$ ;\; b)~$\ens{\dfrac{\pi}{3}\,;\dfrac{2\pi}{3}}$ ;\;
c)~$\ens{-\dfrac{2\pi}{3}\,;\dfrac{2\pi}{3}}$ ;\; d)~$\ens{-\dfrac{3\pi}{4}\,;-\dfrac{\pi}{4}}$ ;\;
e)~$\ens{0}$ ;\; f)~aucune solution.}
\end{methodebox}

\begin{methodebox}[Méthode 10 --- Résoudre $\cos(x)=a$ ou $\sin(x)=a$ sur un autre intervalle]
\etapeexemples
Résoudre, à l'aide du cercle complet :
\begin{questions}[label=\alph*)]
  \item $\cos(x)=\dfrac{\sqrt{2}}{2}$ sur $\intff{0}{2\pi}$ ;
  \item $\sin(x)=-\dfrac{\sqrt{3}}{2}$ sur $\intff{0}{2\pi}$ ;
  \item $\cos(x)=\dfrac{1}{2}$ sur $\intff{-2\pi}{2\pi}$.
\end{questions}
\begin{correction}
\cas{a)}\par\nopagebreak
\begin{minipage}[c]{0.68\linewidth}\raggedright
\explique{On place $\dfrac{\sqrt{2}}{2}$ sur l'axe des abscisses et on trace la droite verticale.}{le cosinus se lit sur l'axe des abscisses}
\explique{Elle coupe le cercle en deux points, de mesures principales $\dfrac{\pi}{4}$ et $-\dfrac{\pi}{4}$.}{lecture sur le cercle complet}
\explique{$\dfrac{\pi}{4}$ est dans $\intff{0}{2\pi}$, mais pas $-\dfrac{\pi}{4}$.}{on compare à l'intervalle demandé}
\begin{calculs}
-\dfrac{\pi}{4}+\dfrac{8\pi}{4} &= \dfrac{7\pi}{4} &\quad& \rem{on ajoute un tour : même point}
\end{calculs}
\explique{$S=\ens{\dfrac{\pi}{4}\,;\dfrac{7\pi}{4}}$.}{une solution par point, dans l'intervalle demandé}
\end{minipage}\hfill
\begin{minipage}[c]{0.28\linewidth}\centering
\equacos{45}{$\frac{\pi}{4}$}{$\frac{7\pi}{4}$}{$\frac{\sqrt{2}}{2}$}
\end{minipage}
\cas{b)}\par\nopagebreak
\begin{minipage}[c]{0.68\linewidth}\raggedright
\explique{On place $-\dfrac{\sqrt{3}}{2}$ sur l'axe des ordonnées et on trace la droite horizontale.}{le sinus se lit sur l'axe des ordonnées}
\explique{Elle coupe le cercle en deux points, de mesures principales $-\dfrac{\pi}{3}$ et $-\dfrac{2\pi}{3}$.}{lecture sur le cercle complet}
\explique{Aucune n'est dans $\intff{0}{2\pi}$.}{elles sont négatives : on ajoute un tour à chacune}
\begin{calculs}
-\dfrac{\pi}{3}+\dfrac{6\pi}{3} &= \dfrac{5\pi}{3} &\quad& \rem{même point, autre mesure}\\
-\dfrac{2\pi}{3}+\dfrac{6\pi}{3} &= \dfrac{4\pi}{3} && \rem{même point, autre mesure}
\end{calculs}
\explique{$S=\ens{\dfrac{4\pi}{3}\,;\dfrac{5\pi}{3}}$.}{on range les solutions}
\end{minipage}\hfill
\begin{minipage}[c]{0.28\linewidth}\centering
\equasin{-60}{-120}{$\frac{5\pi}{3}$}{$\frac{4\pi}{3}$}{$-\frac{\sqrt{3}}{2}$}
\end{minipage}
\cas{c)}\par\nopagebreak
\begin{minipage}[c]{0.68\linewidth}\raggedright
\explique{On place $\dfrac{1}{2}$ sur l'axe des abscisses et on trace la droite verticale.}{le cosinus se lit sur l'axe des abscisses}
\explique{Elle coupe le cercle en deux points, de mesures principales $\dfrac{\pi}{3}$ et $-\dfrac{\pi}{3}$.}{lecture sur le cercle complet}
\explique{Les deux sont dans $\intff{-2\pi}{2\pi}$, qui couvre deux tours : chaque point a une seconde mesure.}{deux tours : deux solutions par point}
\begin{calculs}
\dfrac{\pi}{3}-\dfrac{6\pi}{3} &= -\dfrac{5\pi}{3} &\quad& \rem{un tour de moins}\\
-\dfrac{\pi}{3}+\dfrac{6\pi}{3} &= \dfrac{5\pi}{3} && \rem{un tour de plus}
\end{calculs}
\explique{$S=\ens{-\dfrac{5\pi}{3}\,;-\dfrac{\pi}{3}\,;\dfrac{\pi}{3}\,;\dfrac{5\pi}{3}}$.}{quatre solutions}
\end{minipage}\hfill
\begin{minipage}[c]{0.28\linewidth}\centering
\equacos{60}{$\frac{\pi}{3}$\\$-\frac{5\pi}{3}$}{$-\frac{\pi}{3}$\\$\frac{5\pi}{3}$}{$\frac{1}{2}$}
\end{minipage}
\end{correction}

\etape{À vous de jouer}
Résoudre, à l'aide du cercle complet :
\begin{questions}[label=\alph*)]
  \item $\cos(x)=\dfrac{\sqrt{3}}{2}$ sur $\intff{0}{2\pi}$ ;
  \item $\sin(x)=-\dfrac{1}{2}$ sur $\intff{0}{2\pi}$ ;
  \item $\cos(x)=\dfrac{\sqrt{2}}{2}$ sur $\intff{-2\pi}{2\pi}$.
\end{questions}
\reponse{a)~$\ens{\dfrac{\pi}{6}\,;\dfrac{11\pi}{6}}$ ;\; b)~$\ens{\dfrac{7\pi}{6}\,;\dfrac{11\pi}{6}}$ ;\;
c)~$\ens{-\dfrac{7\pi}{4}\,;-\dfrac{\pi}{4}\,;\dfrac{\pi}{4}\,;\dfrac{7\pi}{4}}$.}
\end{methodebox}

\begin{methodebox}[Méthode 11 --- Retrouver la mesure principale d'après le cosinus et le sinus]
\etapeexemples
Déterminer la mesure principale de l'angle orienté de mesure $x$ tel que :
\begin{questions}[label=\alph*)]
  \item $\cos(x)=-\dfrac{1}{2}$ et $\sin(x)=\dfrac{\sqrt{3}}{2}$ ;
  \item $\cos(x)=\dfrac{\sqrt{2}}{2}$ et $\sin(x)=-\dfrac{\sqrt{2}}{2}$ ;
  \item $\cos(x)=0$ et $\sin(x)=-1$.
\end{questions}
\begin{correction}
\cas{a)}\par\nopagebreak
\begin{minipage}[c]{0.68\linewidth}\raggedright
\explique{On place $-\dfrac{1}{2}$ sur l'axe des abscisses et $\dfrac{\sqrt{3}}{2}$ sur l'axe des ordonnées.}{on trace les deux droites}
\explique{Elles se croisent sur le cercle, au point associé à $\dfrac{2\pi}{3}$.}{un seul point a ces deux coordonnées}
\explique{$x=\dfrac{2\pi}{3}$.}{lecture sur le cercle complet}
\end{minipage}\hfill
\begin{minipage}[c]{0.28\linewidth}\centering
\equapt{120}{$\frac{2\pi}{3}$}{$-\frac{1}{2}$}{$\frac{\sqrt{3}}{2}$}
\end{minipage}
\cas{b)}\par\nopagebreak
\begin{minipage}[c]{0.68\linewidth}\raggedright
\explique{On place $\dfrac{\sqrt{2}}{2}$ sur l'axe des abscisses et $-\dfrac{\sqrt{2}}{2}$ sur l'axe des ordonnées.}{on trace les deux droites}
\explique{Elles se croisent sur le cercle, au point associé à $-\dfrac{\pi}{4}$.}{un seul point a ces deux coordonnées}
\explique{$x=-\dfrac{\pi}{4}$.}{lecture sur le cercle complet}
\end{minipage}\hfill
\begin{minipage}[c]{0.28\linewidth}\centering
\equapt{-45}{$-\frac{\pi}{4}$}{$\frac{\sqrt{2}}{2}$}{$-\frac{\sqrt{2}}{2}$}
\end{minipage}
\cas{c)}\par\nopagebreak
\begin{minipage}[c]{0.68\linewidth}\raggedright
\explique{Le point a pour coordonnées $(0\,;-1)$.}{il est tout en bas du cercle}
\explique{$x=-\dfrac{\pi}{2}$.}{lecture sur le cercle complet}
\end{minipage}\hfill
\begin{minipage}[c]{0.28\linewidth}\centering
\cerclept{-90}{$-\frac{\pi}{2}\qquad$}{}
\end{minipage}
\end{correction}

\etape{À vous de jouer}
Déterminer la mesure principale de l'angle orienté de mesure $x$ tel que :
\begin{questions}[label=\alph*)]
  \item $\cos(x)=-\dfrac{\sqrt{3}}{2}$ et $\sin(x)=\dfrac{1}{2}$ ;
  \item $\cos(x)=\dfrac{1}{2}$ et $\sin(x)=-\dfrac{\sqrt{3}}{2}$ ;
  \item $\cos(x)=-1$ et $\sin(x)=0$.
\end{questions}
\reponse{a)~$\dfrac{5\pi}{6}$ ;\; b)~$-\dfrac{\pi}{3}$ ;\; c)~$\pi$.}
\end{methodebox}

\partiecours{Fonctions cosinus et sinus}

\begin{definitionbox}[Définition 6 --- Fonctions cosinus et sinus, période]
La \textbf{fonction cosinus} associe à tout nombre réel $x$ le nombre $\cos(x)$ ; la
\textbf{fonction sinus} lui associe $\sin(x)$. Leurs courbes s'appellent des
\textbf{sinusoïdes}.

\smallskip
Une fonction $f$ est \textbf{périodique de période $T$} si $f(x+T)=f(x)$ pour tout
nombre réel $x$ : sa courbe reproduit le même morceau sur chaque intervalle de longueur $T$.
\end{definitionbox}

\begin{proprietebox}[Propriété 6 --- Période et variations]
Chaque colonne du tableau de valeurs de la partie V, qui vient du cercle trigonométrique,
donne un point de la courbe (les croix) : par exemple, $\cos\left(\dfrac{\pi}{3}\right)=\dfrac{1}{2}$
donne le point de coordonnées $\left(\dfrac{\pi}{3}\,;\dfrac{1}{2}\right)$ sur la courbe du cosinus.

\smallskip
Comme $\cos(x+2\pi)=\cos(x)$ et $\sin(x+2\pi)=\sin(x)$, les fonctions cosinus et sinus sont
périodiques de période $2\pi$ : le morceau de $\intff{-\pi}{\pi}$ se répète. Dans les repères
ci-dessous, un carreau vaut $\dfrac{\pi}{6}$ en abscisse et $\dfrac{1}{2}$ en ordonnée.
\begin{center}
\begin{tikzpicture}[x=1.2cm,y=1.2cm]
  \reperetrigo{-12}{12}\courbecos{-12}{12}\gradcos\ptstable{cos}
  \node[coulCourbe,font=\footnotesize,fill=white,inner sep=1pt] at (2.2,1.25) {$y=\cos(x)$};
\end{tikzpicture}

\medskip
\begin{tikzpicture}[x=1.2cm,y=1.2cm]
  \reperetrigo{-12}{12}\courbesin{-12}{12}\gradsin\ptstable{sin}
  \node[coulCourbe,font=\footnotesize,fill=white,inner sep=1pt] at (3.3,1.25) {$y=\sin(x)$};
\end{tikzpicture}
\end{center}
Le tableau de valeurs donne aussi le sens de variation : quand $x$ va de $-\pi$ à $0$,
$\cos(x)$ monte de $-1$ à $1$, puis redescend à $-1$ quand $x$ va de $0$ à $\pi$.
Tableaux de variation sur une période, $\intff{-\pi}{\pi}$ :

\medskip
\begin{minipage}[t]{0.48\linewidth}\centering
\begin{tikzpicture}
  \tkzTabInit[lgt=1.5,espcl=2.2]{$x$/1,$\cos(x)$/2}{$-\pi$,$0$,$\pi$}
  \tkzTabVar{-/$-1$,+/$1$,-/$-1$}
\end{tikzpicture}
\end{minipage}\hfill
\begin{minipage}[t]{0.48\linewidth}\centering
\begin{tikzpicture}
  \tkzTabInit[lgt=1.5,espcl=1.5]{$x$/1.3,$\sin(x)$/2}{$-\pi$,$-\dfrac{\pi}{2}$,$\dfrac{\pi}{2}$,$\pi$}
  \tkzTabVar{+/$0$,-/$-1$,+/$1$,-/$0$}
\end{tikzpicture}
\end{minipage}
\end{proprietebox}

\begin{methodebox}[Méthode 12 --- Dresser le tableau de variation sur un autre intervalle]
\etapeexemples
Dresser le tableau de variation, à l'aide des courbes de la propriété 6 :
\begin{questions}[label=\alph*)]
  \item de la fonction cosinus sur $\intff{0}{2\pi}$ ;
  \item de la fonction sinus sur $\intff{-\dfrac{\pi}{2}}{\dfrac{3\pi}{2}}$ ;
  \item de la fonction cosinus sur $\intff{2\pi}{4\pi}$.
\end{questions}
\begin{correction}
\cas{a)}\par\nopagebreak
\begin{minipage}[c]{0.46\linewidth}\raggedright
La courbe part de $1$, descend à $-1$ en $x=\pi$, puis remonte à $1$.
\end{minipage}\hfill
\begin{minipage}[c]{0.50\linewidth}\centering
\begin{tikzpicture}
  \tkzTabInit[lgt=1.5,espcl=2]{$x$/1,$\cos(x)$/2}{$0$,$\pi$,$2\pi$}
  \tkzTabVar{+/$1$,-/$-1$,+/$1$}
\end{tikzpicture}
\end{minipage}
\cas{b)}\par\nopagebreak
\begin{minipage}[c]{0.46\linewidth}\raggedright
La courbe part de $-1$, monte à $1$ en $x=\dfrac{\pi}{2}$, puis redescend à $-1$.
\end{minipage}\hfill
\begin{minipage}[c]{0.50\linewidth}\centering
\begin{tikzpicture}
  \tkzTabInit[lgt=1.5,espcl=2]{$x$/1.3,$\sin(x)$/2}{$-\dfrac{\pi}{2}$,$\dfrac{\pi}{2}$,$\dfrac{3\pi}{2}$}
  \tkzTabVar{-/$-1$,+/$1$,-/$-1$}
\end{tikzpicture}
\end{minipage}
\cas{c)}\par\nopagebreak
\begin{minipage}[c]{0.46\linewidth}\raggedright
La période est $2\pi$ : la courbe refait sur $\intff{2\pi}{4\pi}$ ce qu'elle fait sur
$\intff{0}{2\pi}$. On ajoute $2\pi$ aux abscisses du tableau a).
\end{minipage}\hfill
\begin{minipage}[c]{0.50\linewidth}\centering
\begin{tikzpicture}
  \tkzTabInit[lgt=1.5,espcl=2]{$x$/1,$\cos(x)$/2}{$2\pi$,$3\pi$,$4\pi$}
  \tkzTabVar{+/$1$,-/$-1$,+/$1$}
\end{tikzpicture}
\end{minipage}
\end{correction}

\etape{À vous de jouer}
Dresser le tableau de variation :
\begin{questions}[label=\alph*)]
  \item de la fonction cosinus sur $\intff{-2\pi}{0}$ ;
  \item de la fonction sinus sur $\intff{\dfrac{\pi}{2}}{\dfrac{5\pi}{2}}$ ;
  \item de la fonction cosinus sur $\intff{4\pi}{6\pi}$.
\end{questions}
\reponse{a)~décroissante de $1$ à $-1$ sur $\intff{-2\pi}{-\pi}$, croissante de $-1$ à $1$ sur $\intff{-\pi}{0}$ ;\;
b)~décroissante de $1$ à $-1$ sur $\intff{\dfrac{\pi}{2}}{\dfrac{3\pi}{2}}$, croissante de $-1$ à $1$ sur $\intff{\dfrac{3\pi}{2}}{\dfrac{5\pi}{2}}$ ;\;
c)~décroissante de $1$ à $-1$ sur $\intff{4\pi}{5\pi}$, croissante de $-1$ à $1$ sur $\intff{5\pi}{6\pi}$.}
\end{methodebox}

\begin{methodebox}[Méthode 13 --- Lire la période d'une fonction sur sa courbe]
\etapeexemples
Lire la période de chaque fonction $f$ et $g$, dont voici les courbes.

\smallskip
\begin{minipage}[c]{0.48\linewidth}\centering
\begin{tikzpicture}[x=0.6cm,y=0.6cm]
  \repere{-5}{-3}{5}{3}
  \gradx{-4,-2,2,4}\grady{-2,2}
  \motif{-8,-4,0,4}{(0,0) .. controls (0.35,0.7) and (0.65,2) .. (1,2)
    .. controls (1.35,2) and (1.65,0.7) .. (2,0)
    .. controls (2.35,-0.7) and (2.65,-2) .. (3,-2)
    .. controls (3.35,-2) and (3.65,-0.7) .. (4,0)}
  \node[coulCourbe,font=\footnotesize,fill=white,inner sep=1pt] at (-4,2.5) {$\Cf$};
\end{tikzpicture}
\end{minipage}\hfill
\begin{minipage}[c]{0.48\linewidth}\centering
\begin{tikzpicture}[x=0.6cm,y=0.6cm]
  \repere{-5}{-3}{5}{3}
  \gradx{-4,-2,2,4}\grady{-2,2}
  \motif{-6,-4,-2,0,2,4}{(0,1) .. controls (0.35,1) and (0.65,-1) .. (1,-1)
    .. controls (1.35,-1) and (1.65,1) .. (2,1)}
  \node[coulCourbe,font=\footnotesize,fill=white,inner sep=1pt] at (-4,2.5) {$\Cg$};
\end{tikzpicture}
\end{minipage}
\begin{correction}
\cas{Fonction $f$}
\explique{La période de $f$ est $4$.}{le même morceau revient tous les $4$ carreaux : par exemple, $\Cf$ monte en passant par $0$ en $x=0$, puis en $x=4$}
\cas{Fonction $g$}
\explique{La période de $g$ est $2$.}{les maximums sont en $x=0$, $x=2$, $x=4$ : ils reviennent tous les $2$ carreaux}
\end{correction}

\etape{À vous de jouer}
Lire la période de chaque fonction $h$ et $k$, dont voici les courbes.

\smallskip
\begin{minipage}[c]{0.48\linewidth}\centering
\begin{tikzpicture}[x=0.6cm,y=0.6cm]
  \repere{-5}{-3}{5}{3}
  \gradx{-3,3}\grady{-2,2}
  \motif{-6,-3,0,3}{(0,-1) .. controls (0.35,-1) and (0.65,2) .. (1,2)
    .. controls (1.7,2) and (2.3,-1) .. (3,-1)}
  \node[coulCourbe,font=\footnotesize,fill=white,inner sep=1pt] at (-4,2.5) {$\mathcal{C}_{h}$};
\end{tikzpicture}
\end{minipage}\hfill
\begin{minipage}[c]{0.48\linewidth}\centering
\begin{tikzpicture}[x=0.6cm,y=0.6cm]
  \repere{-5}{-3}{5}{3}
  \gradx{-4,-2,2,4}\grady{-2,2}
  \motif{-8,-4,0,4}{(0,2) .. controls (0.35,2) and (0.65,-1) .. (1,-1)
    .. controls (2,-1) and (3,2) .. (4,2)}
  \node[coulCourbe,font=\footnotesize,fill=white,inner sep=1pt] at (-3.5,-2.5) {$\mathcal{C}_{k}$};
\end{tikzpicture}
\end{minipage}
\reponse{La période de $h$ est $3$ ;\; la période de $k$ est $4$.}
\end{methodebox}

\partiecours{Résolution graphique d'équations et d'inéquations}

Comme pour toute fonction, on lit les solutions sur la courbe : les abscisses de ses points
d'intersection avec une droite horizontale. La courbe montre d'un seul coup toutes les
solutions d'un intervalle.

\begin{methodebox}[Méthode 14 --- Résoudre graphiquement une équation]
\etapeexemples
Un carreau vaut $\dfrac{\pi}{6}$ en abscisse et $\dfrac{1}{2}$ en ordonnée. Résoudre graphiquement :
\begin{questions}[label=\alph*)]
  \item $\cos(x)=\dfrac{1}{2}$ sur $\intff{-2\pi}{2\pi}$ ;
  \item $\sin(x)=-\dfrac{1}{2}$ sur $\intff{-2\pi}{2\pi}$ ;
  \item $\sin(x)=\dfrac{\sqrt{3}}{2}$ sur $\intff{0}{2\pi}$.
\end{questions}
\begin{center}
\begin{tikzpicture}[x=1.2cm,y=1.2cm]
  \reperetrigo{-12}{12}\courbecos{-12}{12}\gradcos
  \node[coulCourbe,font=\footnotesize,fill=white,inner sep=1pt] at (2.2,1.25) {$y=\cos(x)$};
\end{tikzpicture}

\medskip
\begin{tikzpicture}[x=1.2cm,y=1.2cm]
  \reperetrigo{-12}{12}\courbesin{-12}{12}\gradsin
  \node[coulCourbe,font=\footnotesize,fill=white,inner sep=1pt] at (3.3,1.25) {$y=\sin(x)$};
\end{tikzpicture}
\end{center}
\begin{correction}
\cas{a)}\par\nopagebreak
\explique{On trace la droite d'équation $y=\dfrac{1}{2}$ et on lit les abscisses des points d'intersection.}{la droite coupe la courbe en quatre points}
\explique{$S=\ens{-\dfrac{5\pi}{3}\,;-\dfrac{\pi}{3}\,;\dfrac{\pi}{3}\,;\dfrac{5\pi}{3}}$.}{$\tfrac{\pi}{3}$ : deux carreaux à droite de $0$}
\begin{center}
\begin{tikzpicture}[x=1.2cm,y=1.2cm]
  \reperetrigo{-12}{12}\courbecos{-12}{12}
  \droiteh{-12}{12}{0.5}{$y=\frac{1}{2}$}
  \lirex{-10}{0.5}\lirex{-2}{0.5}\lirex{2}{0.5}\lirex{10}{0.5}
  \gradpi{-12/-2\pi,-10/-\frac{5\pi}{3},-2/-\frac{\pi}{3},2/\frac{\pi}{3},10/\frac{5\pi}{3},12/2\pi}
\end{tikzpicture}
\end{center}
\cas{b)}\par\nopagebreak
\explique{On trace la droite d'équation $y=-\dfrac{1}{2}$ et on lit les abscisses des points d'intersection.}{même démarche, sur la courbe du sinus}
\explique{$S=\ens{-\dfrac{5\pi}{6}\,;-\dfrac{\pi}{6}\,;\dfrac{7\pi}{6}\,;\dfrac{11\pi}{6}}$.}{$\tfrac{7\pi}{6}$ : sept carreaux à droite de $0$}
\begin{center}
\begin{tikzpicture}[x=1.2cm,y=1.2cm]
  \reperetrigo{-12}{12}\courbesin{-12}{12}
  \droiteh{-12}{12}{-0.5}{$y=-\frac{1}{2}$}
  \lirex{-5}{-0.5}\lirex{-1}{-0.5}\lirex{7}{-0.5}\lirex{11}{-0.5}
  \gradpi[above]{-12/-2\pi,-5/-\frac{5\pi}{6},-1/-\frac{\pi}{6},7/\frac{7\pi}{6},11/\frac{11\pi}{6}}
\end{tikzpicture}
\end{center}
\cas{c)}\par\nopagebreak
\begin{minipage}[c]{0.42\linewidth}\raggedright
\explique{$S=\ens{\dfrac{\pi}{3}\,;\dfrac{2\pi}{3}}$.}{seul $\intff{0}{2\pi}$ compte}
\end{minipage}\hfill
\begin{minipage}[c]{0.56\linewidth}\centering
\begin{tikzpicture}[x=1.2cm,y=1.2cm]
  \reperetrigo{0}{12}\courbesin{0}{12}
  \droiteh{0}{12}{0.866}{$y=\frac{\sqrt{3}}{2}$}
  \lirex{2}{0.866}\lirex{4}{0.866}
  \gradpi{2/\frac{\pi}{3},4/\frac{2\pi}{3},12/2\pi}
\end{tikzpicture}
\end{minipage}
\end{correction}

\etape{À vous de jouer}
Résoudre graphiquement, à l'aide des courbes ci-dessus :
\begin{center}
a)~$\cos(x)=-\dfrac{1}{2}$ sur $\intff{-2\pi}{2\pi}$ \qquad
b)~$\sin(x)=\dfrac{1}{2}$ sur $\intff{-2\pi}{2\pi}$ \qquad
c)~$\cos(x)=\dfrac{\sqrt{3}}{2}$ sur $\intff{0}{2\pi}$
\end{center}
\reponse{a)~$\ens{-\dfrac{4\pi}{3}\,;-\dfrac{2\pi}{3}\,;\dfrac{2\pi}{3}\,;\dfrac{4\pi}{3}}$ ;\;
b)~$\ens{-\dfrac{11\pi}{6}\,;-\dfrac{7\pi}{6}\,;\dfrac{\pi}{6}\,;\dfrac{5\pi}{6}}$ ;\;
c)~$\ens{\dfrac{\pi}{6}\,;\dfrac{11\pi}{6}}$.}
\end{methodebox}

\begin{methodebox}[Méthode 15 --- Résoudre graphiquement une inéquation]
\etapeexemples
Résoudre graphiquement, à l'aide des courbes de la méthode 14 :
\begin{questions}[label=\alph*)]
  \item $\cos(x)\ge\dfrac{1}{2}$ sur $\intff{-\pi}{\pi}$ ;
  \item $\sin(x)<\dfrac{1}{2}$ sur $\intff{0}{2\pi}$ ;
  \item $\cos(x)\le-\dfrac{1}{2}$ sur $\intff{-2\pi}{2\pi}$.
\end{questions}
\begin{correction}
\cas{a)}\par\nopagebreak
\begin{minipage}[c]{0.42\linewidth}\raggedright
\explique{$S=\intff{-\dfrac{\pi}{3}}{\dfrac{\pi}{3}}$.}{la courbe est sur ou au-dessus de la droite ; crochets fermés à cause de $\ge$}
\end{minipage}\hfill
\begin{minipage}[c]{0.56\linewidth}\centering
\begin{tikzpicture}[x=1.2cm,y=1.2cm]
  \reperetrigo{-6}{6}\courbecos{-6}{6}
  \droiteh{-6}{6}{0.5}{$y=\frac{1}{2}$}
  \interaxe{-2}{2}
  \lirex{-2}{0.5}\lirex{2}{0.5}
  \gradpi{-6/-\pi,-2/-\frac{\pi}{3},2/\frac{\pi}{3},6/\pi}
\end{tikzpicture}
\end{minipage}
\cas{b)}\par\nopagebreak
\begin{minipage}[c]{0.42\linewidth}\raggedright
\explique{$S=\intfo{0}{\dfrac{\pi}{6}}\cup\intof{\dfrac{5\pi}{6}}{2\pi}$.}{la courbe est au-dessous de la droite, en deux morceaux ; crochets ouverts à cause de $<$}
\end{minipage}\hfill
\begin{minipage}[c]{0.56\linewidth}\centering
\begin{tikzpicture}[x=1.2cm,y=1.2cm]
  \reperetrigo{0}{12}\courbesin{0}{12}
  \droiteh{0}{12}{0.5}{$y=\frac{1}{2}$}
  \interaxe{0}{1}\interaxe[coulCourbeC]{5}{12}
  \lirex{1}{0.5}\lirex{5}{0.5}
  \gradpi{1/\frac{\pi}{6},5/\frac{5\pi}{6},12/2\pi}
\end{tikzpicture}
\end{minipage}
\cas{c)}\par\nopagebreak
\explique{$S=\intff{-\dfrac{4\pi}{3}}{-\dfrac{2\pi}{3}}\cup\intff{\dfrac{2\pi}{3}}{\dfrac{4\pi}{3}}$.}{la courbe est sur ou au-dessous de la droite}
\begin{center}
\begin{tikzpicture}[x=1.2cm,y=1.2cm]
  \reperetrigo{-12}{12}\courbecos{-12}{12}
  \droiteh{-12}{12}{-0.5}{$y=-\frac{1}{2}$}
  \interaxe{-8}{-4}\interaxe[coulCourbeC]{4}{8}
  \lirex{-8}{-0.5}\lirex{-4}{-0.5}\lirex{4}{-0.5}\lirex{8}{-0.5}
  \gradpi[above]{-8/-\frac{4\pi}{3},-4/-\frac{2\pi}{3},4/\frac{2\pi}{3},8/\frac{4\pi}{3}}
\end{tikzpicture}
\end{center}
\end{correction}

\etape{À vous de jouer}
Résoudre graphiquement, à l'aide des courbes de la méthode 14 :
\begin{center}
a)~$\cos(x)\ge-\dfrac{1}{2}$ sur $\intff{-\pi}{\pi}$ \qquad
b)~$\sin(x)>-\dfrac{1}{2}$ sur $\intff{0}{2\pi}$ \qquad
c)~$\sin(x)\ge\dfrac{1}{2}$ sur $\intff{-2\pi}{2\pi}$
\end{center}
\reponse{a)~$\intff{-\dfrac{2\pi}{3}}{\dfrac{2\pi}{3}}$ ;\;
b)~$\intfo{0}{\dfrac{7\pi}{6}}\cup\intof{\dfrac{11\pi}{6}}{2\pi}$ ;\;
c)~$\intff{-\dfrac{11\pi}{6}}{-\dfrac{7\pi}{6}}\cup\intff{\dfrac{\pi}{6}}{\dfrac{5\pi}{6}}$.}
\end{methodebox}

\partiecours{Fonctions paires et impaires}

\begin{definitionbox}[Définition 7 --- Fonction paire, fonction impaire]
\begin{minipage}[t]{0.48\linewidth}\raggedright
Une fonction $f$ est \textbf{paire} si, pour tout nombre réel $x$ :
\[ f(-x)=f(x) \]
Sa courbe est symétrique par rapport à l'axe des ordonnées.
\end{minipage}\hfill
\begin{minipage}[t]{0.48\linewidth}\raggedright
Une fonction $f$ est \textbf{impaire} si, pour tout nombre réel $x$ :
\[ f(-x)=-f(x) \]
Sa courbe est symétrique par rapport à l'origine du repère.
\end{minipage}
\end{definitionbox}

\begin{proprietebox}[Propriété 7 --- Parité du cosinus et du sinus]
D'après les angles associés (propriété 5), $\cos(-x)=\cos(x)$ et $\sin(-x)=-\sin(x)$ :
la fonction cosinus est \textbf{paire} et la fonction sinus est \textbf{impaire}. On le voit sur
leurs courbes (propriété 6).
\end{proprietebox}

\begin{methodebox}[Méthode 16 --- Reconnaître sur la courbe une fonction paire ou impaire]
\etapeexemples
Dire si chaque fonction $f$, $g$ et $h$ est paire, impaire ou ni l'une ni l'autre.

\smallskip
\begin{minipage}[c]{0.32\linewidth}\centering
\begin{tikzpicture}[x=0.5cm,y=0.5cm]
  \repere{-4}{-3}{4}{3}
  \gradx{-3,-1,1,3}\grady{-2,2}
  \draw[coulCourbe,line width=1.1pt] (-3,-2) .. controls (-2.4,-1) and (-1.6,2) .. (-1,2)
    .. controls (-0.6,2) and (-0.4,1) .. (0,1)
    .. controls (0.4,1) and (0.6,2) .. (1,2)
    .. controls (1.6,2) and (2.4,-1) .. (3,-2);
  \node[coulCourbe,font=\footnotesize,fill=white,inner sep=1pt] at (2.5,2.3) {$\Cf$};
\end{tikzpicture}
\end{minipage}\hfill
\begin{minipage}[c]{0.32\linewidth}\centering
\begin{tikzpicture}[x=0.5cm,y=0.5cm]
  \repere{-4}{-3}{4}{3}
  \gradx{-3,-1,1,3}\grady{-2,2}
  \draw[coulCourbe,line width=1.1pt] (-3,1) .. controls (-2.3,0.4) and (-1.7,-2) .. (-1,-2)
    .. controls (-0.4,-2) and (0.4,2) .. (1,2)
    .. controls (1.7,2) and (2.3,-0.4) .. (3,-1);
  \node[coulCourbe,font=\footnotesize,fill=white,inner sep=1pt] at (2.5,2.3) {$\Cg$};
\end{tikzpicture}
\end{minipage}\hfill
\begin{minipage}[c]{0.32\linewidth}\centering
\begin{tikzpicture}[x=0.5cm,y=0.5cm]
  \repere{-4}{-3}{4}{3}
  \gradx{-3,-1,1,3}\grady{-2,2}
  \draw[coulCourbe,line width=1.1pt] (-3,-2) .. controls (-2,-1) and (-1,2) .. (0,2)
    .. controls (1,2) and (2,-0.4) .. (3,-1);
  \node[coulCourbe,font=\footnotesize,fill=white,inner sep=1pt] at (2.5,2.3) {$\mathcal{C}_{h}$};
\end{tikzpicture}
\end{minipage}
\begin{correction}
\cas{Fonction $f$}
\explique{La fonction $f$ est paire.}{plier le long de l'axe des ordonnées : les deux moitiés se superposent, par exemple $f(-1)$ et $f(1)$ valent $2$}
\cas{Fonction $g$}
\explique{La fonction $g$ est impaire.}{demi-tour autour de l'origine : $g(1)$ vaut $2$ et $g(-1)$ vaut $-2$}
\cas{Fonction $h$}
\explique{La fonction $h$ n'est ni paire ni impaire.}{$h(-3)$ vaut $-2$ et $h(3)$ vaut $-1$ : ni égaux ni opposés, donc aucune symétrie}
\end{correction}

\etape{À vous de jouer}
Dire si chaque fonction $k$, $m$ et $n$ est paire, impaire ou ni l'une ni l'autre.

\smallskip
\begin{minipage}[c]{0.32\linewidth}\centering
\begin{tikzpicture}[x=0.5cm,y=0.5cm]
  \repere{-4}{-3}{4}{3}
  \gradx{-3,-1,1,3}\grady{-2,2}
  \draw[coulCourbe,line width=1.1pt] (-3,3) .. controls (-2.2,1) and (-1,-1) .. (0,-1)
    .. controls (1,-1) and (2.2,1) .. (3,3);
  \node[coulCourbe,font=\footnotesize,fill=white,inner sep=1pt] at (1,2.4) {$\mathcal{C}_{k}$};
\end{tikzpicture}
\end{minipage}\hfill
\begin{minipage}[c]{0.32\linewidth}\centering
\begin{tikzpicture}[x=0.5cm,y=0.5cm]
  \repere{-4}{-3}{4}{3}
  \gradx{-3,-1,1,3}\grady{-2,2}
  \draw[coulCourbe,line width=1.1pt] (-3,-2) .. controls (-2.4,-1) and (-1.6,1) .. (-1,1)
    .. controls (-0.5,1) and (0.5,-1) .. (1,-1)
    .. controls (1.6,-1) and (2.4,1) .. (3,2);
  \node[coulCourbe,font=\footnotesize,fill=white,inner sep=1pt] at (2.4,2.6) {$\mathcal{C}_{m}$};
\end{tikzpicture}
\end{minipage}\hfill
\begin{minipage}[c]{0.32\linewidth}\centering
\begin{tikzpicture}[x=0.5cm,y=0.5cm]
  \repere{-4}{-3}{4}{3}
  \gradx{-3,-1,1,2,3}\grady{-2,2}
  \draw[coulCourbe,line width=1.1pt] (-3,1) .. controls (-2.4,1.6) and (-1.6,3) .. (-1,3)
    .. controls (0,3) and (1,-2) .. (2,-2)
    .. controls (2.4,-2) and (2.7,-0.6) .. (3,0);
  \node[coulCourbe,font=\footnotesize,fill=white,inner sep=1pt] at (2.5,2.3) {$\mathcal{C}_{n}$};
\end{tikzpicture}
\end{minipage}
\reponse{$k$ est paire ;\; $m$ est impaire ;\; $n$ n'est ni paire ni impaire.}
\end{methodebox}

\begin{methodebox}[Méthode 17 --- Démontrer qu'une fonction est paire ou impaire]
\textbf{But :} montrer que $f(-x)=f(x)$ si $f$ est paire, ou que $f(-x)=-f(x)$ si $f$ est impaire.
Pour cela, on remplace $x$ par $-x$ dans la fonction.
\etapeexemples
Ces fonctions sont définies sur $\R$.
\begin{questions}[label=\alph*)]
  \item Démontrer que $f(x)=x^{2}+3\cos(x)$ est paire.
  \item Démontrer que $g(x)=2x-\sin(x)$ est impaire.
  \item Démontrer que $h(x)=\sin(3x)-x$ est impaire.
\end{questions}
\begin{correction}
\cas{a)}
\begin{calculs}
f(-x) &= (-x)^{2}+3\cos(-x) &\qquad& \rem{on remplace $x$ par $(-x)$}\\
f(-x) &= x^{2}+3\cos(x) && \rem{$(-x)^{2}=x^{2}$ et $\cos(-x)=\cos(x)$}\\
f(-x) &= f(x) && \rem{on reconnaît $f(x)$}
\end{calculs}
\explique{La fonction $f$ est paire.}{on conclut par une phrase}
\cas{b)}
\begin{calculs}
g(-x) &= 2\times(-x)-\sin(-x) &\qquad& \rem{on remplace $x$ par $(-x)$}\\
g(-x) &= -2x+\sin(x) && \rem{$\sin(-x)=-\sin(x)$}\\
g(-x) &= -\big(2x-\sin(x)\big) && \rem{on met le signe $-$ en facteur}\\
g(-x) &= -g(x) && \rem{on reconnaît $g(x)$}
\end{calculs}
\explique{La fonction $g$ est impaire.}{on conclut par une phrase}
\cas{c)}
\begin{calculs}
h(-x) &= \sin\big(3\times(-x)\big)-(-x) &\qquad& \rem{on remplace $x$ par $(-x)$}\\
h(-x) &= \sin(-3x)+x && \rem{on simplifie}\\
h(-x) &= -\sin(3x)+x && \rem{$\sin(-X)=-\sin(X)$, avec $X=3x$}\\
h(-x) &= -\big(\sin(3x)-x\big) && \rem{on met le signe $-$ en facteur}\\
h(-x) &= -h(x) && \rem{on reconnaît $h(x)$}
\end{calculs}
\explique{La fonction $h$ est impaire.}{on conclut par une phrase}
\end{correction}

\etape{À vous de jouer}
Ces fonctions sont définies sur $\R$.
\begin{questions}[label=\alph*)]
  \item Démontrer que $k(x)=x^{2}-2\cos(x)$ est paire.
  \item Démontrer que $m(x)=3x+\sin(x)$ est impaire.
  \item Démontrer que $n(x)=\sin(2x)+x$ est impaire.
\end{questions}
\reponse{a)~$k(-x)=k(x)$ ;\; b)~$m(-x)=-m(x)$ ;\; c)~$n(-x)=-n(x)$.}
\end{methodebox}

\begin{methodebox}[Méthode 18 --- Compléter une courbe par parité et périodicité]
\etapeexemples
\begin{minipage}[c]{0.50\linewidth}\raggedright
La fonction $f$ est paire et périodique de période $4$. Voici sa courbe sur $\intff{0}{2}$.
Compléter la courbe sur $\intff{-4}{4}$.
\end{minipage}\hfill
\begin{minipage}[c]{0.46\linewidth}\centering
\begin{tikzpicture}[x=0.55cm,y=0.55cm]
  \repere{-5}{-3}{5}{3}
  \gradx{-4,-2,2,4}\grady{-2,2}
  \draw[coulCourbe,line width=1.1pt] (0,2) .. controls (0.7,2) and (1.3,-2) .. (2,-2);
\end{tikzpicture}
\end{minipage}
\begin{correction}
\cas{Résolution}\par\nopagebreak
\begin{minipage}[c]{0.50\linewidth}\raggedright
\explique{Sur $\intff{-2}{0}$ : on trace le symétrique par rapport à l'axe des ordonnées.}{$f$ est paire (en rouge)}
\explique{Sur $\intff{2}{4}$ et $\intff{-4}{-2}$ : on reproduit le morceau de $\intff{-2}{2}$, décalé de $4$.}{la période est $4$ (en vert)}
\end{minipage}\hfill
\begin{minipage}[c]{0.46\linewidth}\centering
\begin{tikzpicture}[x=0.55cm,y=0.55cm]
  \repere{-5}{-3}{5}{3}
  \gradx{-4,-2,2,4}\grady{-2,2}
  \draw[coulCourbe,line width=1.1pt] (0,2) .. controls (0.7,2) and (1.3,-2) .. (2,-2);
  \draw[coulCourbeB,line width=1.1pt] (0,2) .. controls (-0.7,2) and (-1.3,-2) .. (-2,-2);
  \draw[coulCourbeC,line width=1.1pt] (2,-2) .. controls (2.7,-2) and (3.3,2) .. (4,2);
  \draw[coulCourbeC,line width=1.1pt] (-2,-2) .. controls (-2.7,-2) and (-3.3,2) .. (-4,2);
\end{tikzpicture}
\end{minipage}
\end{correction}

\etape{À vous de jouer}
\begin{minipage}[c]{0.50\linewidth}\raggedright
La fonction $g$ est impaire et périodique de période $4$. Voici sa courbe sur $\intff{0}{2}$.
Compléter la courbe sur $\intff{-4}{4}$.
\end{minipage}\hfill
\begin{minipage}[c]{0.46\linewidth}\centering
\begin{tikzpicture}[x=0.55cm,y=0.55cm]
  \repere{-5}{-3}{5}{3}
  \gradx{-4,-2,2,4}\grady{-2,2}
  \draw[coulCourbe,line width=1.1pt] (0,0) .. controls (0.35,0.7) and (0.65,2) .. (1,2)
    .. controls (1.35,2) and (1.65,0.7) .. (2,0);
\end{tikzpicture}
\end{minipage}
\reponse{La courbe passe par $(-3\,;2)$, $(-2\,;0)$, $(-1\,;-2)$, $(0\,;0)$, $(1\,;2)$, $(2\,;0)$, $(3\,;-2)$ et $(4\,;0)$.}
\end{methodebox}

\partiecours{Équations du type $\cos(3x)=a$}

Dans $\cos(2x)$, l'angle est $2x$. On résout d'abord sur le cercle, pour l'angle tout entier, qu'on nomme
$X$, puis on revient à $x$. Attention : $X$ ne parcourt pas le même intervalle que $x$.

\begin{methodebox}[Méthode 19 --- Résoudre $\cos(kx)=a$ ou $\sin(kx)=a$ sur un intervalle]
\etapeexemples
Résoudre, à l'aide du cercle complet :
\begin{questions}[label=\alph*)]
  \item $\cos(2x)=\dfrac{1}{2}$ sur $\intff{0}{\pi}$ ;
  \item $\sin(3x)=1$ sur $\intff{-\pi}{\pi}$ ;
  \item $\cos\left(\dfrac{x}{2}\right)=0$ sur $\intff{0}{4\pi}$.
\end{questions}
\begin{correction}
\cas{a)}\par\nopagebreak
\begin{minipage}[c]{0.68\linewidth}\raggedright
\explique{On pose $X=2x$.}{on nomme l'angle}
\explique{$x$ va de $0$ à $\pi$, donc $X$ va de $0$ à $2\pi$.}{on multiplie les bornes par $2$}
\explique{$\cos(X)=\dfrac{1}{2}$ donne $X=\dfrac{\pi}{3}$ ou $X=\dfrac{5\pi}{3}$.}{sur le cercle : la droite verticale passant par $\tfrac{1}{2}$ ; $-\tfrac{\pi}{3}$ s'écrit $\tfrac{5\pi}{3}$ sur $\intff{0}{2\pi}$}
\end{minipage}\hfill
\begin{minipage}[c]{0.28\linewidth}\centering
\equacos{60}{$\frac{\pi}{3}$}{$\frac{5\pi}{3}$}{$\frac{1}{2}$}
\end{minipage}
\begin{calculs}
2x &= \dfrac{\pi}{3} &\qquad& \rem{premier cas : $X$ vaut $2x$}\\
x &= \dfrac{\pi}{6} && \rem{on divise par $2$}\\
2x &= \dfrac{5\pi}{3} && \rem{second cas}\\
x &= \dfrac{5\pi}{6} && \rem{on divise par $2$}
\end{calculs}
\explique{$S=\ens{\dfrac{\pi}{6}\,;\dfrac{5\pi}{6}}$.}{les deux valeurs sont bien dans $\intff{0}{\pi}$}
\cas{b)}\par\nopagebreak
\begin{minipage}[c]{0.68\linewidth}\raggedright
\explique{On pose $X=3x$.}{on nomme l'angle}
\explique{$x$ va de $-\pi$ à $\pi$, donc $X$ va de $-3\pi$ à $3\pi$.}{on multiplie les bornes par $3$}
\explique{$\sin(X)=1$ donne $X=-\dfrac{3\pi}{2}$, $X=\dfrac{\pi}{2}$ ou $X=\dfrac{5\pi}{2}$.}{sur le cercle : le point tout en haut, $\tfrac{\pi}{2}$ ; un tour de moins, un tour de plus restent dans l'intervalle de $X$}
\end{minipage}\hfill
\begin{minipage}[c]{0.28\linewidth}\centering
\cerclept{90}{$\frac{\pi}{2}\qquad$}{\horizontale{90}}
\end{minipage}
\begin{calculs}
3x &= -\dfrac{3\pi}{2} &\qquad& \rem{premier cas}\\
x &= -\dfrac{\pi}{2} && \rem{on divise par $3$}\\
3x &= \dfrac{\pi}{2} && \rem{deuxième cas}\\
x &= \dfrac{\pi}{6} && \rem{on divise par $3$}\\
3x &= \dfrac{5\pi}{2} && \rem{troisième cas}\\
x &= \dfrac{5\pi}{6} && \rem{on divise par $3$}
\end{calculs}
\explique{$S=\ens{-\dfrac{\pi}{2}\,;\dfrac{\pi}{6}\,;\dfrac{5\pi}{6}}$.}{trois solutions : $X$ parcourt trois tours}
\cas{c)}\par\nopagebreak
\begin{minipage}[c]{0.68\linewidth}\raggedright
\explique{On pose $X=\dfrac{x}{2}$.}{on nomme l'angle}
\explique{$x$ va de $0$ à $4\pi$, donc $X$ va de $0$ à $2\pi$.}{on divise les bornes par $2$}
\explique{$\cos(X)=0$ donne $X=\dfrac{\pi}{2}$ ou $X=\dfrac{3\pi}{2}$.}{sur le cercle : les points tout en haut et tout en bas ; $-\tfrac{\pi}{2}$ s'écrit $\tfrac{3\pi}{2}$ sur $\intff{0}{2\pi}$}
\end{minipage}\hfill
\begin{minipage}[c]{0.28\linewidth}\centering
\begin{tikzpicture}[x=1.25cm,y=1.25cm]
  \cerclebase\ptc{90}{$\frac{\pi}{2}\qquad$}\ptc{-90}{$\frac{3\pi}{2}\qquad$}
\end{tikzpicture}
\end{minipage}
\begin{calculs}
\dfrac{x}{2} &= \dfrac{\pi}{2} &\qquad& \rem{premier cas}\\
x &= \pi && \rem{on multiplie par $2$}\\
\dfrac{x}{2} &= \dfrac{3\pi}{2} && \rem{second cas}\\
x &= 3\pi && \rem{on multiplie par $2$}
\end{calculs}
\explique{$S=\ens{\pi\,;3\pi}$.}{les deux valeurs sont bien dans $\intff{0}{4\pi}$}
\end{correction}

\etape{À vous de jouer}
Résoudre, à l'aide du cercle complet :
\begin{questions}[label=\alph*)]
  \item $\cos(2x)=\dfrac{\sqrt{3}}{2}$ sur $\intff{0}{\pi}$ ;
  \item $\sin(3x)=-1$ sur $\intff{-\pi}{\pi}$ ;
  \item $\cos\left(\dfrac{x}{2}\right)=\dfrac{1}{2}$ sur $\intff{0}{4\pi}$.
\end{questions}
\reponse{a)~$\ens{\dfrac{\pi}{12}\,;\dfrac{11\pi}{12}}$ ;\; b)~$\ens{-\dfrac{5\pi}{6}\,;-\dfrac{\pi}{6}\,;\dfrac{\pi}{2}}$ ;\; c)~$\ens{\dfrac{2\pi}{3}\,;\dfrac{10\pi}{3}}$.}
\end{methodebox}

\begin{methodebox}[Méthode 20 --- Résoudre une équation dont l'angle contient un terme en plus]
\etapeexemples
Résoudre, à l'aide du cercle complet :
\begin{questions}[label=\alph*)]
  \item $\sin\left(x+\dfrac{\pi}{6}\right)=1$ sur $\intff{-\pi}{\pi}$ ;
  \item $\sin\left(2x-\dfrac{\pi}{2}\right)=\dfrac{1}{2}$ sur $\intff{0}{\pi}$.
\end{questions}
\begin{correction}
\cas{a)}\par\nopagebreak
\begin{minipage}[c]{0.68\linewidth}\raggedright
\explique{On pose $X=x+\dfrac{\pi}{6}$.}{on nomme l'angle}
\explique{$X$ va de $-\dfrac{5\pi}{6}$ à $\dfrac{7\pi}{6}$.}{on ajoute $\tfrac{\pi}{6}$ aux bornes $-\pi$ et $\pi$}
\explique{$\sin(X)=1$ donne $X=\dfrac{\pi}{2}$.}{sur le cercle : le point tout en haut, seul dans cet intervalle}
\end{minipage}\hfill
\begin{minipage}[c]{0.28\linewidth}\centering
\cerclept{90}{$\frac{\pi}{2}\qquad$}{\horizontale{90}}
\end{minipage}
\begin{calculs}
x+\dfrac{\pi}{6} &= \dfrac{\pi}{2} &\qquad& \rem{on revient à $x$}\\
x &= \dfrac{3\pi}{6}-\dfrac{\pi}{6} && \rem{on retire $\tfrac{\pi}{6}$ de chaque côté}\\
x &= \dfrac{\pi}{3} && \rem{$\tfrac{2\pi}{6}$ se simplifie}
\end{calculs}
\explique{$S=\ens{\dfrac{\pi}{3}}$.}{la valeur est bien dans $\intff{-\pi}{\pi}$}
\cas{b)}\par\nopagebreak
\begin{minipage}[c]{0.68\linewidth}\raggedright
\explique{On pose $X=2x-\dfrac{\pi}{2}$.}{on nomme l'angle}
\explique{$X$ va de $-\dfrac{\pi}{2}$ à $\dfrac{3\pi}{2}$.}{on remplace $x$ par les bornes $0$ et $\pi$}
\explique{$\sin(X)=\dfrac{1}{2}$ donne $X=\dfrac{\pi}{6}$ ou $X=\dfrac{5\pi}{6}$.}{sur le cercle : la droite horizontale passant par $\tfrac{1}{2}$}
\end{minipage}\hfill
\begin{minipage}[c]{0.28\linewidth}\centering
\equasin{30}{150}{$\frac{\pi}{6}$}{$\frac{5\pi}{6}$}{$\frac{1}{2}$}
\end{minipage}
\begin{calculs}
2x-\dfrac{\pi}{2} &= \dfrac{\pi}{6} &\qquad& \rem{premier cas}\\
2x &= \dfrac{\pi}{6}+\dfrac{3\pi}{6} && \rem{on ajoute $\tfrac{\pi}{2}$, soit $\tfrac{3\pi}{6}$}\\
2x &= \dfrac{2\pi}{3} && \rem{$\tfrac{4\pi}{6}$ se simplifie}\\
x &= \dfrac{\pi}{3} && \rem{on divise par $2$}\\
2x-\dfrac{\pi}{2} &= \dfrac{5\pi}{6} && \rem{second cas}\\
2x &= \dfrac{5\pi}{6}+\dfrac{3\pi}{6} && \rem{on ajoute $\tfrac{3\pi}{6}$}\\
2x &= \dfrac{4\pi}{3} && \rem{$\tfrac{8\pi}{6}$ se simplifie}\\
x &= \dfrac{2\pi}{3} && \rem{on divise par $2$}
\end{calculs}
\explique{$S=\ens{\dfrac{\pi}{3}\,;\dfrac{2\pi}{3}}$.}{les deux valeurs sont bien dans $\intff{0}{\pi}$}
\end{correction}

\etape{À vous de jouer}
Résoudre, à l'aide du cercle complet :
\begin{questions}[label=\alph*)]
  \item $\cos\left(x-\dfrac{\pi}{6}\right)=1$ sur $\intff{-\pi}{\pi}$ ;
  \item $\sin\left(2x-\dfrac{\pi}{2}\right)=\dfrac{\sqrt{3}}{2}$ sur $\intff{0}{\pi}$.
\end{questions}
\reponse{a)~$\ens{\dfrac{\pi}{6}}$ ;\; b)~$\ens{\dfrac{5\pi}{12}\,;\dfrac{7\pi}{12}}$.}
\end{methodebox}

\partiecours{Fonctions sinusoïdales : amplitude, période, phase}

La tension d'une prise de courant, le son d'un haut-parleur, la vibration d'un moteur :
ces grandeurs oscillent au cours du temps et se décrivent toutes par un cosinus ou un sinus.

\begin{definitionbox}[Définition 8 --- Fonction sinusoïdale]
\begin{minipage}[c]{0.58\linewidth}\raggedright
Une \textbf{fonction sinusoïdale} du temps $t$ s'écrit
\begin{center}
$A\cos(\omega t+\varphi)$ \quad ou \quad $A\sin(\omega t+\varphi)$
\end{center}
où $A$ et $\omega$ sont des nombres positifs et $\varphi$ un nombre réel.
\begin{itemize}
  \item $A$ est l'\textbf{amplitude} : la valeur maximale.
  \item $\omega$ est la \textbf{pulsation}, en rad/s.
  \item $\omega t+\varphi$ est la \textbf{phase instantanée}.
  \item $\varphi$ est la \textbf{phase à l'origine}, en rad.
\end{itemize}
\end{minipage}\hfill
\begin{minipage}[c]{0.38\linewidth}\centering
\begin{tikzpicture}[x=0.6cm,y=0.6cm]
  \reperet{-1}{-3}{7}{3}
  \gradx{1,5}\grady{2}
  \courbet{-1}{7}{2*sin(90*\x)}
  \draw[red,line width=1.4pt] (1,0)--(1,2);
  \node[red,font=\footnotesize,left,inner sep=2pt] at (1,1.1) {$A$};
  \draw[{Stealth[length=1.8mm]}-{Stealth[length=1.8mm]},coulCourbeC,line width=0.9pt] (1,2.5)--(5,2.5);
  \node[coulCourbeC,font=\footnotesize,above,inner sep=1.5pt,fill=white] at (3,2.5) {$T$};
\end{tikzpicture}
\end{minipage}
\end{definitionbox}

\begin{proprietebox}[Propriété 8 --- Période et fréquence]
Une fonction sinusoïdale de pulsation $\omega$ a pour \textbf{période} $T=\dfrac{2\pi}{\omega}$.
Si $t$ est en secondes, sa \textbf{fréquence}, en hertz (Hz), vaut $\dfrac{1}{T}$ : c'est le
nombre de périodes par seconde.

\smallskip
\demo{Quand $t$ augmente de $T$, la phase augmente d'un tour :}
\begin{calculs}
\omega(t+T)+\varphi &= \omega t+\omega T+\varphi &\qquad& \rem{on développe}\\
\omega(t+T)+\varphi &= \omega t+\omega\times\dfrac{2\pi}{\omega}+\varphi && \rem{on remplace $T$}\\
\omega(t+T)+\varphi &= \omega t+\varphi+2\pi && \rem{un tour de plus : même cosinus, même sinus}
\end{calculs}
\end{proprietebox}

\begin{methodebox}[Méthode 21 --- Lire l'amplitude et la phase, calculer la période]
\etapeexemples
Pour chaque fonction, donner l'amplitude, la pulsation et la phase à l'origine, puis calculer
la période :
\begin{questions}[label=\alph*)]
  \item la tension d'une prise de courant, en volts, $t$ en secondes : $u(t)=325\cos(100\pi t)$ (donner aussi la fréquence) ;
  \item $v(t)=3\sin\left(2t+\dfrac{\pi}{4}\right)$ ;
  \item $w(t)=5\cos\left(\dfrac{\pi}{2}t-\dfrac{\pi}{3}\right)$.
\end{questions}
\begin{correction}
\cas{a)}
\explique{L'amplitude est $325$ V, la pulsation $100\pi$ rad/s, la phase à l'origine $0$ rad.}{le nombre devant le cosinus, celui devant $t$ ; aucun terme ajouté}
\begin{calculs}
T &= \dfrac{2\pi}{\omega} &\qquad& \rem{la formule}\\
T &= \dfrac{2\barre{\pi}}{100\barre{\pi}} && \rem{on remplace $\omega$ par $100\pi$}\\
T &= 0{,}02\text{ s} && \rem{$2$ divisé par $100$}\\
\dfrac{1}{T} &= 50\text{ Hz} && \rem{la fréquence : $1$ divisé par $0{,}02$}
\end{calculs}
\cas{b)}
\explique{L'amplitude est $3$, la pulsation $2$, la phase à l'origine $\dfrac{\pi}{4}$.}{on lit $A$, $\omega$ et $\varphi$ dans l'expression}
\begin{calculs}
T &= \dfrac{2\pi}{\omega} &\qquad& \rem{la formule}\\
T &= \dfrac{2\pi}{2} && \rem{on remplace $\omega$ par $2$}\\
T &= \pi &&
\end{calculs}
\cas{c)}
\explique{L'amplitude est $5$, la pulsation $\dfrac{\pi}{2}$, la phase à l'origine $-\dfrac{\pi}{3}$.}{le signe $-$ fait partie de la phase}
\begin{calculs}
T &= \dfrac{2\pi}{\omega} &\qquad& \rem{la formule}\\
T &= 2\pi\times\dfrac{2}{\pi} && \rem{diviser par $\tfrac{\pi}{2}$, c'est multiplier par $\tfrac{2}{\pi}$}\\
T &= 4 && \rem{les $\pi$ se simplifient}
\end{calculs}
\end{correction}

\etape{À vous de jouer}
Pour chaque fonction, donner l'amplitude, la pulsation et la phase à l'origine, puis calculer
la période :
\begin{questions}[label=\alph*)]
  \item la tension du réseau électrique américain, en volts, $t$ en secondes : $u(t)=170\cos(120\pi t)$ (donner aussi la fréquence) ;
  \item $v(t)=4\sin\left(3t+\dfrac{\pi}{6}\right)$ ;
  \item $w(t)=2\cos\left(\dfrac{\pi}{4}t-\dfrac{\pi}{2}\right)$.
\end{questions}
\reponse{a)~$170$ V ; $120\pi$ rad/s ; $0$ rad ; $T=\dfrac{1}{60}$ s ; $60$ Hz ;\;
b)~$4$ ; $3$ ; $\dfrac{\pi}{6}$ ; $T=\dfrac{2\pi}{3}$ ;\; c)~$2$ ; $\dfrac{\pi}{4}$ ; $-\dfrac{\pi}{2}$ ; $T=8$.}
\end{methodebox}

\partiecours{Courbe d'une fonction sinusoïdale}

La courbe d'une fonction sinusoïdale oscille entre $-A$ et $A$ et se répète à chaque période.
La phase à l'origine la décale le long de l'axe du temps.

\begin{methodebox}[Méthode 22 --- Tracer la courbe d'une fonction sinusoïdale]
\etapeexemples
Pour chaque fonction $g$ et $h$, calculer les images de $0$, $1$, $2$, $3$ et $4$, puis tracer sa
courbe sur $\intff{-4}{4}$ :
\begin{center}
$g(t)=2\sin\left(\dfrac{\pi}{2}t\right)$ \qquad $h(t)=2\cos\left(\dfrac{\pi}{2}t+\dfrac{\pi}{2}\right)$
\end{center}
\begin{correction}
\cas{Fonction $g$}\par\nopagebreak
\begin{minipage}[c]{0.62\linewidth}\raggedright
\begin{calculs}
g(0) &= 2\sin(0) = 0 &\quad& \rem{$\tfrac{\pi}{2}\times 0$ vaut $0$}\\
g(1) &= 2\sin\left(\dfrac{\pi}{2}\right) = 2 && \rem{$\tfrac{\pi}{2}\times 1$ vaut $\tfrac{\pi}{2}$}\\
g(2) &= 2\sin(\pi) = 0 && \rem{$\tfrac{\pi}{2}\times 2$ vaut $\pi$}\\
g(3) &= 2\sin\left(\dfrac{3\pi}{2}\right) = -2 && \rem{le bas du cercle : sinus $-1$}\\
g(4) &= 2\sin(2\pi) = 0 && \rem{un tour complet}
\end{calculs}
\explique{La période est $4$ : on reproduit le morceau de $\intff{0}{4}$ sur $\intff{-4}{0}$.}{$\tfrac{2\pi}{\omega}$, avec $\omega=\tfrac{\pi}{2}$, vaut $4$}
\end{minipage}\hfill
\begin{minipage}[c]{0.34\linewidth}\centering
\sinusoidepts{2*sin(90*\x)}{-2,2}{0/0,1/2,2/0,3/-2,4/0}
\end{minipage}
\cas{Fonction $h$}\par\nopagebreak
\begin{minipage}[c]{0.62\linewidth}\raggedright
\begin{calculs}
h(0) &= 2\cos\left(\dfrac{\pi}{2}\right) = 0 &\quad& \rem{la phase à l'origine seule}\\
h(1) &= 2\cos(\pi) = -2 && \rem{$\tfrac{\pi}{2}+\tfrac{\pi}{2}$ vaut $\pi$}\\
h(2) &= 2\cos\left(\dfrac{3\pi}{2}\right) = 0 && \rem{$\pi+\tfrac{\pi}{2}$ vaut $\tfrac{3\pi}{2}$}\\
h(3) &= 2\cos(2\pi) = 2 && \rem{$\tfrac{3\pi}{2}+\tfrac{\pi}{2}$ vaut $2\pi$}\\
h(4) &= 2\cos\left(\dfrac{5\pi}{2}\right) = 0 && \rem{un tour et $\tfrac{\pi}{2}$}
\end{calculs}
\explique{Même amplitude et même période que $g$ : la phase décale la courbe.}{on compare les deux figures}
\end{minipage}\hfill
\begin{minipage}[c]{0.34\linewidth}\centering
\sinusoidepts{2*cos(90*\x+90)}{-2,2}{0/0,1/-2,2/0,3/2,4/0}
\end{minipage}
\end{correction}

\etape{À vous de jouer}
Pour chaque fonction $k$ et $m$, calculer les images de $0$, $1$, $2$, $3$ et $4$, puis tracer sa
courbe sur $\intff{-4}{4}$ :
\begin{center}
$k(t)=3\sin\left(\dfrac{\pi}{2}t\right)$ \qquad $m(t)=3\cos\left(\dfrac{\pi}{2}t+\dfrac{\pi}{2}\right)$
\end{center}
\reponse{Images de $k$ : $0$ ; $3$ ; $0$ ; $-3$ ; $0$ ;\; images de $m$ : $0$ ; $-3$ ; $0$ ; $3$ ; $0$.}
\end{methodebox}

\begin{methodebox}[Méthode 23 --- Retrouver l'expression d'une fonction sinusoïdale sur sa courbe]
\etapeexemples
Retrouver l'expression de chaque fonction, avec $\varphi$ dans $\intof{-\pi}{\pi}$ :

\smallskip
\begin{minipage}[t]{0.32\linewidth}\centering
a)~$f(t)=A\cos(\omega t+\varphi)$\par\smallskip
\sinusoide{3*cos(90*\x)}{-3,3}
\end{minipage}\hfill
\begin{minipage}[t]{0.32\linewidth}\centering
b)~$g(t)=A\sin(\omega t+\varphi)$\par\smallskip
\sinusoide{2*sin(180*\x+90)}{-2,2}
\end{minipage}\hfill
\begin{minipage}[t]{0.32\linewidth}\centering
c)~$h(t)=A\cos(\omega t+\varphi)$\par\smallskip
\sinusoide{2*cos(90*\x+90)}{-2,2}
\end{minipage}
\begin{correction}
\cas{a)}
\explique{$A=3$ et $T=4$.}{le maximum ; deux maximums de suite sont en $t=0$ et $t=4$}
\begin{calculs}
\omega &= \dfrac{2\pi}{T} &\qquad& \rem{la formule de la période, à l'envers}\\
\omega &= \dfrac{\pi}{2} && \rem{$\tfrac{2\pi}{4}$ se simplifie}\\
f(0) &= 3\cos(\varphi) && \rem{en $t=0$, il reste la phase}\\
3\cos(\varphi) &= 3 && \rem{on lit $f(0)=3$ sur la courbe}\\
\cos(\varphi) &= 1 && \rem{on divise par $3$}\\
\varphi &= 0 && \rem{le point $(1\,;0)$}
\end{calculs}
\explique{$f(t)=3\cos\left(\dfrac{\pi}{2}t\right)$.}{on remplace $A$, $\omega$ et $\varphi$}
\cas{b)}
\explique{$A=2$ et $T=2$.}{deux maximums de suite sont en $t=0$ et $t=2$}
\begin{calculs}
\omega &= \dfrac{2\pi}{T} &\qquad& \rem{la formule}\\
\omega &= \pi && \rem{$\tfrac{2\pi}{2}$ se simplifie}\\
2\sin(\varphi) &= 2 && \rem{on lit $g(0)=2$ sur la courbe}\\
\sin(\varphi) &= 1 && \rem{on divise par $2$}\\
\varphi &= \dfrac{\pi}{2} && \rem{le point tout en haut du cercle}
\end{calculs}
\explique{$g(t)=2\sin\left(\pi t+\dfrac{\pi}{2}\right)$.}{on remplace $A$, $\omega$ et $\varphi$}
\cas{c)}
\explique{$A=2$ et $T=4$, donc $\omega=\dfrac{\pi}{2}$.}{deux maximums de suite sont en $t=-1$ et $t=3$}
\begin{calculs}
2\cos(\varphi) &= 0 &\qquad& \rem{on lit $h(0)=0$ sur la courbe}\\
\cos(\varphi) &= 0 && \rem{deux valeurs possibles : $-\tfrac{\pi}{2}$ et $\tfrac{\pi}{2}$}\\
h(1) &= 2\cos\left(\dfrac{\pi}{2}+\varphi\right) && \rem{un second point pour choisir}\\
2\cos\left(\dfrac{\pi}{2}+\varphi\right) &= -2 && \rem{on lit $h(1)=-2$ sur la courbe}\\
\cos\left(\dfrac{\pi}{2}+\varphi\right) &= -1 && \rem{donc $\tfrac{\pi}{2}+\varphi$ vaut $\pi$}\\
\varphi &= \dfrac{\pi}{2} && \rem{$-\tfrac{\pi}{2}$ donnerait $\cos(0)=1$ : ne convient pas}
\end{calculs}
\explique{$h(t)=2\cos\left(\dfrac{\pi}{2}t+\dfrac{\pi}{2}\right)$.}{on remplace $A$, $\omega$ et $\varphi$}
\end{correction}

\etape{À vous de jouer}
Retrouver l'expression de chaque fonction, avec $\varphi$ dans $\intof{-\pi}{\pi}$ :

\smallskip
\begin{minipage}[t]{0.32\linewidth}\centering
a)~$k(t)=A\cos(\omega t+\varphi)$\par\smallskip
\sinusoide{2*cos(180*\x)}{-2,2}
\end{minipage}\hfill
\begin{minipage}[t]{0.32\linewidth}\centering
b)~$m(t)=A\sin(\omega t+\varphi)$\par\smallskip
\sinusoide{3*sin(90*\x+90)}{-3,3}
\end{minipage}\hfill
\begin{minipage}[t]{0.32\linewidth}\centering
c)~$n(t)=A\cos(\omega t+\varphi)$\par\smallskip
\sinusoide{3*cos(90*\x-90)}{-3,3}
\end{minipage}
\reponse{a)~$k(t)=2\cos(\pi t)$ ;\; b)~$m(t)=3\sin\left(\dfrac{\pi}{2}t+\dfrac{\pi}{2}\right)$ ;\;
c)~$n(t)=3\cos\left(\dfrac{\pi}{2}t-\dfrac{\pi}{2}\right)$.}
\end{methodebox}

\end{document}
